
\Data\ID{1003102501}
\Updated{2020-07-18 01:07:40}
\Level{A}
\SubArea{Quadratic equations with complex roots}
\LANGen{
\Question{
Find the set of all complex roots of the following quadratic equation.
\[ 9x^2 + 2 = 0 \]}
{\( \left\{-\frac{\sqrt2}3\mathrm{i}; \frac{\sqrt2}3\mathrm{i}\right\} \)}{\( \left\{-\sqrt{\frac{2}3}\mathrm{i}; \sqrt{\frac23}\mathrm{i}\right\} \)}{\( \left\{-\frac23\mathrm{i}; \frac23\mathrm{i}\right\} \)}{\( \left\{-\frac29\mathrm{i}; \frac29\mathrm{i}\right\} \)}{\( \emptyset \)}{}}
\LANGpl{
\Question{
Wskaż zbiór wszystkich pierwiastków zespolonych danego równania kwadratowego.
\[ 9x^2 + 2 = 0 \]}
{\( \left\{-\frac{\sqrt2}3\mathrm{i}; \frac{\sqrt2}3\mathrm{i}\right\} \)}{\( \left\{-\sqrt{\frac{2}3}\mathrm{i}; \sqrt{\frac23}\mathrm{i}\right\} \)}{\( \left\{-\frac23\mathrm{i}; \frac23\mathrm{i}\right\} \)}{\( \left\{-\frac29\mathrm{i}; \frac29\mathrm{i}\right\} \)}{\( \emptyset \)}{}}
\LANGcs{
\Question{
Určete množinu komplexních kořenů dané rovnice.
\[ 9x^2 + 2 = 0 \]}
{\( \left\{-\frac{\sqrt2}3\mathrm{i}; \frac{\sqrt2}3\mathrm{i}\right\} \)}{\( \left\{-\sqrt{\frac{2}3}\mathrm{i}; \sqrt{\frac23}\mathrm{i}\right\} \)}{\( \left\{-\frac23\mathrm{i}; \frac23\mathrm{i}\right\} \)}{\( \left\{-\frac29\mathrm{i}; \frac29\mathrm{i}\right\} \)}{\( \emptyset \)}{}}
\LANGsk{
\Question{
Určte množinu komplexných koreňov danej rovnice.
\[ 9x^2 + 2 = 0 \]}
{\( \left\{-\frac{\sqrt2}3\mathrm{i}; \frac{\sqrt2}3\mathrm{i}\right\} \)}{\( \left\{-\sqrt{\frac{2}3}\mathrm{i}; \sqrt{\frac23}\mathrm{i}\right\} \)}{\( \left\{-\frac23\mathrm{i}; \frac23\mathrm{i}\right\} \)}{\( \left\{-\frac29\mathrm{i}; \frac29\mathrm{i}\right\} \)}{\( \emptyset \)}{}}
\LANGes{
\Question{
Halla el conjunto de todas las raíces complejas de la siguiente ecuación cuadrática.
\[ 9x^2 + 2 = 0 \]}
{\( \left\{-\frac{\sqrt2}3\mathrm{i}; \frac{\sqrt2}3\mathrm{i}\right\} \)}{\( \left\{-\sqrt{\frac{2}3}\mathrm{i}; \sqrt{\frac23}\mathrm{i}\right\} \)}{\( \left\{-\frac23\mathrm{i}; \frac23\mathrm{i}\right\} \)}{\( \left\{-\frac29\mathrm{i}; \frac29\mathrm{i}\right\} \)}{\( \emptyset \)}{}}
\End

\Data\ID{1003102502}
\Updated{2020-07-18 01:07:46}
\Level{A}
\SubArea{Quadratic equations with complex roots}
\LANGen{
\Question{
Find the complex roots of the following quadratic equation.
\[ 4x^2 + 0.0025 = 0 \]}
{\( x_1=-\frac1{40}\mathrm{i}\text{, } x_2=\frac1{40}\mathrm{i} \)}{\( x_1=-\frac1{4}\mathrm{i}\text{, } x_2=\frac1{4}\mathrm{i} \)}{\( x_1=-\frac5{16}\mathrm{i}\text{, } x_2=\frac5{16}\mathrm{i} \)}{\( x_1=-\frac5{4}\mathrm{i}\text{, } x_2=\frac5{4}\mathrm{i} \)}{}{}}
\LANGpl{
\Question{
Wyznacz pierwiastki zespolone danego równania kwadratowego.
\[ 4x^2 + 0{,}0025 = 0 \]}
{\( x_1=-\frac1{40}\mathrm{i}\text{, } x_2=\frac1{40}\mathrm{i} \)}{\( x_1=-\frac1{4}\mathrm{i}\text{, } x_2=\frac1{4}\mathrm{i} \)}{\( x_1=-\frac5{16}\mathrm{i}\text{, } x_2=\frac5{16}\mathrm{i} \)}{\( x_1=-\frac5{4}\mathrm{i}\text{, } x_2=\frac5{4}\mathrm{i} \)}{}{}}
\LANGcs{
\Question{
Určete komplexní kořeny následující kvadratické rovnice.
\[ 4x^2 + 0{,}0025 = 0 \]}
{\( x_1=-\frac1{40}\mathrm{i}\text{, } x_2=\frac1{40}\mathrm{i} \)}{\( x_1=-\frac1{4}\mathrm{i}\text{, } x_2=\frac1{4}\mathrm{i} \)}{\( x_1=-\frac5{16}\mathrm{i}\text{, } x_2=\frac5{16}\mathrm{i} \)}{\( x_1=-\frac5{4}\mathrm{i}\text{, } x_2=\frac5{4}\mathrm{i} \)}{}{}}
\LANGsk{
\Question{
Určte komplexné korene danej kvadratickej rovnice.
\[ 4x^2 + 0{,}0025 = 0 \]}
{\( x_1=-\frac1{40}\mathrm{i}\text{, } x_2=\frac1{40}\mathrm{i} \)}{\( x_1=-\frac1{4}\mathrm{i}\text{, } x_2=\frac1{4}\mathrm{i} \)}{\( x_1=-\frac5{16}\mathrm{i}\text{, } x_2=\frac5{16}\mathrm{i} \)}{\( x_1=-\frac5{4}\mathrm{i}\text{, } x_2=\frac5{4}\mathrm{i} \)}{}{}}
\LANGes{
\Question{
Halla las raíces complejas de la siguiente ecuación cuadrática.
\[ 4x^2 + 0.0025 = 0 \]}
{\( x_1=-\frac1{40}\mathrm{i}\text{, } x_2=\frac1{40}\mathrm{i} \)}{\( x_1=-\frac1{4}\mathrm{i}\text{, } x_2=\frac1{4}\mathrm{i} \)}{\( x_1=-\frac5{16}\mathrm{i}\text{, } x_2=\frac5{16}\mathrm{i} \)}{\( x_1=-\frac5{4}\mathrm{i}\text{, } x_2=\frac5{4}\mathrm{i} \)}{}{}}
\End

\Data\ID{1003102503}
\Updated{2020-07-18 01:07:04}
\Level{A}
\SubArea{Quadratic equations with complex roots}
\LANGen{
\Question{
Find the complex roots of the following quadratic equation.
\[ 5x^2 + 12 = 0 \]}
{\( x_1=-\frac{2\sqrt{15}}5\mathrm{i}\text{, }x_2=\frac{2\sqrt{15}}5\mathrm{i} \)}{\( x_1=-\frac{\sqrt{15}}5\mathrm{i}\text{, }x_2=\frac{\sqrt{15}}5\mathrm{i} \)}{\( x_1=-\frac{\sqrt{12}}5\mathrm{i}\text{, }x_2=\frac{\sqrt{12}}5\mathrm{i} \)}{\( x_1=-\frac{2\sqrt{3}}5\mathrm{i}\text{, }x_2=\frac{2\sqrt{3}}5\mathrm{i} \)}{}{}}
\LANGpl{
\Question{
Wyznacz pierwiastki zespolone danego równania kwadratowego.
\[ 5x^2 + 12 = 0 \]}
{\( x_1=-\frac{2\sqrt{15}}5\mathrm{i}\text{, }x_2=\frac{2\sqrt{15}}5\mathrm{i} \)}{\( x_1=-\frac{\sqrt{15}}5\mathrm{i}\text{, }x_2=\frac{\sqrt{15}}5\mathrm{i} \)}{\( x_1=-\frac{\sqrt{12}}5\mathrm{i}\text{, }x_2=\frac{\sqrt{12}}5\mathrm{i} \)}{\( x_1=-\frac{2\sqrt{3}}5\mathrm{i}\text{, }x_2=\frac{2\sqrt{3}}5\mathrm{i} \)}{}{}}
\LANGcs{
\Question{
Určete komplexní kořeny následující kvadratické rovnice.
\[ 5x^2 + 12 = 0 \]}
{\( x_1=-\frac{2\sqrt{15}}5\mathrm{i}\text{, }x_2=\frac{2\sqrt{15}}5\mathrm{i} \)}{\( x_1=-\frac{\sqrt{15}}5\mathrm{i}\text{, }x_2=\frac{\sqrt{15}}5\mathrm{i} \)}{\( x_1=-\frac{\sqrt{12}}5\mathrm{i}\text{, }x_2=\frac{\sqrt{12}}5\mathrm{i} \)}{\( x_1=-\frac{2\sqrt{3}}5\mathrm{i}\text{, }x_2=\frac{2\sqrt{3}}5\mathrm{i} \)}{}{}}
\LANGsk{
\Question{
Určte komplexné korene danej kvadratickej rovnice.
\[ 5x^2 + 12 = 0 \]}
{\( x_1=-\frac{2\sqrt{15}}5\mathrm{i}\text{, }x_2=\frac{2\sqrt{15}}5\mathrm{i} \)}{\( x_1=-\frac{\sqrt{15}}5\mathrm{i}\text{, }x_2=\frac{\sqrt{15}}5\mathrm{i} \)}{\( x_1=-\frac{\sqrt{12}}5\mathrm{i}\text{, }x_2=\frac{\sqrt{12}}5\mathrm{i} \)}{\( x_1=-\frac{2\sqrt{3}}5\mathrm{i}\text{, }x_2=\frac{2\sqrt{3}}5\mathrm{i} \)}{}{}}
\LANGes{
\Question{
Halla las raíces complejas de la siguiente ecuación cuadrática.
\[ 5x^2 + 12 = 0 \]}
{\( x_1=-\frac{2\sqrt{15}}5\mathrm{i}\text{, }x_2=\frac{2\sqrt{15}}5\mathrm{i} \)}{\( x_1=-\frac{\sqrt{15}}5\mathrm{i}\text{, }x_2=\frac{\sqrt{15}}5\mathrm{i} \)}{\( x_1=-\frac{\sqrt{12}}5\mathrm{i}\text{, }x_2=\frac{\sqrt{12}}5\mathrm{i} \)}{\( x_1=-\frac{2\sqrt{3}}5\mathrm{i}\text{, }x_2=\frac{2\sqrt{3}}5\mathrm{i} \)}{}{}}
\End

\Data\ID{1003102504}
\Updated{2020-07-18 01:07:26}
\Level{A}
\SubArea{Quadratic equations with complex roots}
\LANGen{
\Question{
Find the set of all complex roots of the following quadratic equation.
\[ x^2 + 4x + 8 = 0 \]}
{\( \left\{ 2\sqrt2\left(\cos\frac{3\pi}4+\mathrm{i}\cdot\sin\frac{3\pi}4\right); 2\sqrt2\left(\cos\frac{5\pi}4+\mathrm{i}\cdot\sin\frac{5\pi}4\right) \right\} \)}{\( \left\{ 2\left(\cos\frac{3\pi}4+\mathrm{i}\cdot\sin\frac{3\pi}4\right); 2\left(\cos\frac{5\pi}4+\mathrm{i}\cdot\sin\frac{5\pi}4\right) \right\} \)}{\( \left\{ 2\sqrt2\left(\cos\frac{\pi}4+\mathrm{i}\cdot\sin\frac{\pi}4\right); 2\sqrt2\left(\cos\frac{7\pi}4+\mathrm{i}\cdot\sin\frac{7\pi}4\right) \right\} \)}{\( \left\{ 2\left(\cos\frac{\pi}4+\mathrm{i}\cdot\sin\frac{\pi}4\right); 2\left(\cos\frac{7\pi}4+\mathrm{i}\cdot\sin\frac{7\pi}4\right) \right\} \)}{}{}}
\LANGpl{
\Question{
Wskaż zbiór wszystkich pierwiastków zespolonych danego równania kwadratowego.
\[ x^2 + 4x + 8 = 0 \]}
{\( \left\{ 2\sqrt2\left(\cos\frac{3\pi}4+\mathrm{i}\cdot\sin\frac{3\pi}4\right); 2\sqrt2\left(\cos\frac{5\pi}4+\mathrm{i}\cdot\sin\frac{5\pi}4\right) \right\} \)}{\( \left\{ 2\left(\cos\frac{3\pi}4+\mathrm{i}\cdot\sin\frac{3\pi}4\right); 2\left(\cos\frac{5\pi}4+\mathrm{i}\cdot\sin\frac{5\pi}4\right) \right\} \)}{\( \left\{ 2\sqrt2\left(\cos\frac{\pi}4+\mathrm{i}\cdot\sin\frac{\pi}4\right); 2\sqrt2\left(\cos\frac{7\pi}4+\mathrm{i}\cdot\sin\frac{7\pi}4\right) \right\} \)}{\( \left\{ 2\left(\cos\frac{\pi}4+\mathrm{i}\cdot\sin\frac{\pi}4\right); 2\left(\cos\frac{7\pi}4+\mathrm{i}\cdot\sin\frac{7\pi}4\right) \right\} \)}{}{}}
\LANGcs{
\Question{
Určete množinu komplexních kořenů dané rovnice.
\[ x^2 + 4x + 8 = 0 \]}
{\( \left\{ 2\sqrt2\left(\cos\frac{3\pi}4+\mathrm{i}\cdot\sin\frac{3\pi}4\right); 2\sqrt2\left(\cos\frac{5\pi}4+\mathrm{i}\cdot\sin\frac{5\pi}4\right) \right\} \)}{\( \left\{ 2\left(\cos\frac{3\pi}4+\mathrm{i}\cdot\sin\frac{3\pi}4\right); 2\left(\cos\frac{5\pi}4+\mathrm{i}\cdot\sin\frac{5\pi}4\right) \right\} \)}{\( \left\{ 2\sqrt2\left(\cos\frac{\pi}4+\mathrm{i}\cdot\sin\frac{\pi}4\right); 2\sqrt2\left(\cos\frac{7\pi}4+\mathrm{i}\cdot\sin\frac{7\pi}4\right) \right\} \)}{\( \left\{ 2\left(\cos\frac{\pi}4+\mathrm{i}\cdot\sin\frac{\pi}4\right); 2\left(\cos\frac{7\pi}4+\mathrm{i}\cdot\sin\frac{7\pi}4\right) \right\} \)}{}{}}
\LANGsk{
\Question{
Určte množinu komplexných koreňov danej rovnice.
\[ x^2 + 4x + 8 = 0 \]}
{\( \left\{ 2\sqrt2\left(\cos\frac{3\pi}4+\mathrm{i}\cdot\sin\frac{3\pi}4\right); 2\sqrt2\left(\cos\frac{5\pi}4+\mathrm{i}\cdot\sin\frac{5\pi}4\right) \right\} \)}{\( \left\{ 2\left(\cos\frac{3\pi}4+\mathrm{i}\cdot\sin\frac{3\pi}4\right); 2\left(\cos\frac{5\pi}4+\mathrm{i}\cdot\sin\frac{5\pi}4\right) \right\} \)}{\( \left\{ 2\sqrt2\left(\cos\frac{\pi}4+\mathrm{i}\cdot\sin\frac{\pi}4\right); 2\sqrt2\left(\cos\frac{7\pi}4+\mathrm{i}\cdot\sin\frac{7\pi}4\right) \right\} \)}{\( \left\{ 2\left(\cos\frac{\pi}4+\mathrm{i}\cdot\sin\frac{\pi}4\right); 2\left(\cos\frac{7\pi}4+\mathrm{i}\cdot\sin\frac{7\pi}4\right) \right\} \)}{}{}}
\LANGes{
\Question{
Halla el conjunto de todas las raíces complejas de la siguiente ecuación cuadrática.
\[ x^2 + 4x + 8 = 0 \]}
{\( \left\{ 2\sqrt2\left(\cos\frac{3\pi}4+\mathrm{i}\cdot\sin\frac{3\pi}4\right); 2\sqrt2\left(\cos\frac{5\pi}4+\mathrm{i}\cdot\sin\frac{5\pi}4\right) \right\} \)}{\( \left\{ 2\left(\cos\frac{3\pi}4+\mathrm{i}\cdot\sin\frac{3\pi}4\right); 2\left(\cos\frac{5\pi}4+\mathrm{i}\cdot\sin\frac{5\pi}4\right) \right\} \)}{\( \left\{ 2\sqrt2\left(\cos\frac{\pi}4+\mathrm{i}\cdot\sin\frac{\pi}4\right); 2\sqrt2\left(\cos\frac{7\pi}4+\mathrm{i}\cdot\sin\frac{7\pi}4\right) \right\} \)}{\( \left\{ 2\left(\cos\frac{\pi}4+\mathrm{i}\cdot\sin\frac{\pi}4\right); 2\left(\cos\frac{7\pi}4+\mathrm{i}\cdot\sin\frac{7\pi}4\right) \right\} \)}{}{}}
\End

\Data\ID{1003102505}
\Updated{2020-07-18 01:07:48}
\Level{A}
\SubArea{Quadratic equations with complex roots}
\LANGen{
\Question{
Find the set of all complex roots of the following equation.
\[ \left(x^2- 2x + 5\right)\left(x^2 + 6x + 10\right) = 0 \]}
{\( \{1\pm2\mathrm{i}; -3\pm\mathrm{i} \} \)}{\( \{-1\pm2\mathrm{i}; -3\pm\mathrm{i} \} \)}{\( \{-1\pm2\mathrm{i}; 3\pm\mathrm{i} \} \)}{\( \{1\pm2\mathrm{i}; 3\pm\mathrm{i} \} \)}{}{}}
\LANGpl{
\Question{
Wskaż zbiór wszystkich pierwiastków zespolonych danego równania kwadratowego.
\[ \left(x^2 - 2x + 5\right)\left(x^2 + 6x + 10\right) = 0 \]}
{\( \{1\pm2\mathrm{i}; -3\pm\mathrm{i} \} \)}{\( \{-1\pm2\mathrm{i}; -3\pm\mathrm{i} \} \)}{\( \{-1\pm2\mathrm{i}; 3\pm\mathrm{i} \} \)}{\( \{1\pm2\mathrm{i}; 3\pm\mathrm{i} \} \)}{}{}}
\LANGcs{
\Question{
Určete množinu všech komplexních kořenů rovnice.
\[ \left(x^2- 2x + 5\right).\left(x^2 + 6x + 10\right) = 0 \]}
{\( \{1\pm2\mathrm{i}; -3\pm\mathrm{i} \} \)}{\( \{-1\pm2\mathrm{i}; -3\pm\mathrm{i} \} \)}{\( \{-1\pm2\mathrm{i}; 3\pm\mathrm{i} \} \)}{\( \{1\pm2\mathrm{i}; 3\pm\mathrm{i} \} \)}{}{}}
\LANGsk{
\Question{
Určte množinu všetkých komplexných koreňov rovnice
\[ \left(x^2 - 2x + 5\right).\left(x^2 + 6x + 10\right) = 0 \]}
{\( \{1\pm2\mathrm{i}; -3\pm\mathrm{i} \} \)}{\( \{-1\pm2\mathrm{i}; -3\pm\mathrm{i} \} \)}{\( \{-1\pm2\mathrm{i}; 3\pm\mathrm{i} \} \)}{\( \{1\pm2\mathrm{i}; 3\pm\mathrm{i} \} \)}{}{}}
\LANGes{
\Question{
Halla el conjunto de todas las raíces complejas de la siguiente ecuación cuadrática.
\[ \left(x^2- 2x + 5\right)\left(x^2 + 6x + 10\right) = 0 \]}
{\( \{1\pm2\mathrm{i}; -3\pm\mathrm{i} \} \)}{\( \{-1\pm2\mathrm{i}; -3\pm\mathrm{i} \} \)}{\( \{-1\pm2\mathrm{i}; 3\pm\mathrm{i} \} \)}{\( \{1\pm2\mathrm{i}; 3\pm\mathrm{i} \} \)}{}{}}
\End

\Data\ID{9000019808}
\Updated{2020-07-18 02:07:30}
\Level{A}
\SubArea{Quadratic equations with complex roots}
\LANGen{
\Question{
Assuming \(x\in \mathbb{C}\),
find the solution set of the following equation. 
\[ 
 x\left (x + 1\right )\left (x^{2} + 1\right ) = 0
\]}
{\(\left \{-1;0;-\mathrm{i};\mathrm{i}\right \}\)}{\(\left \{-1;0;1;-\mathrm{i};\mathrm{i}\right \}\)}{\(\left \{-1;1;-\mathrm{i};\mathrm{i}\right \}\)}{\(\left \{-1;0;-\mathrm{i}\right \}\)}{}{}}
\LANGpl{
\Question{
Zakładając, że \(x\in \mathbb{C}\),
wyznacz zbiór rozwiązań następującego równania.
\[ 
 x\left (x + 1\right )\left (x^{2} + 1\right ) = 0
\]}
{\(\left \{-1;0;-\mathrm{i};\mathrm{i}\right \}\)}{\(\left \{-1;0;1;-\mathrm{i};\mathrm{i}\right \}\)}{\(\left \{-1;1;-\mathrm{i};\mathrm{i}\right \}\)}{\(\left \{-1;0;-\mathrm{i}\right \}\)}{}{}}
\LANGcs{
\Question{
Určete množinu všech řešení rovnice
\(x\left (x + 1\right )\left (x^{2} + 1\right ) = 0\) v
oboru komplexních čísel.}
{\(\left \{-1;0;-\mathrm{i};\mathrm{i}\right \}\)}{\(\left \{-1;0;1;-\mathrm{i};\mathrm{i}\right \}\)}{\(\left \{-1;1;-\mathrm{i};\mathrm{i}\right \}\)}{\(\left \{-1;0;-\mathrm{i}\right \}\)}{}{}}
\LANGsk{
\Question{
Nájdite množinu všetkých riešení danej rovnice pre  \(x\in \mathbb{C}\). 
\[ 
 x\left (x + 1\right )\left (x^{2} + 1\right ) = 0
\]}
{\(\left \{-1;0;-\mathrm{i};\mathrm{i}\right \}\)}{\(\left \{-1;0;1;-\mathrm{i};\mathrm{i}\right \}\)}{\(\left \{-1;1;-\mathrm{i};\mathrm{i}\right \}\)}{\(\left \{-1;0;-\mathrm{i}\right \}\)}{}{}}
\LANGes{
\Question{
Determina el conjunto de todas las soluciones de la ecuación
\(x\left (x + 1\right )\left (x^{2} + 1\right ) = 0\) en el conjunto de los números complejos.}
{\(\left \{-1;0;-\mathrm{i};\mathrm{i}\right \}\)}{\(\left \{-1;0;1;-\mathrm{i};\mathrm{i}\right \}\)}{\(\left \{-1;1;-\mathrm{i};\mathrm{i}\right \}\)}{\(\left \{-1;0;-\mathrm{i}\right \}\)}{}{}}
\End

\Data\ID{9000022803}
\Updated{2020-07-18 02:07:34}
\Level{A}
\SubArea{Quadratic equations with complex roots}
\LANGen{
\Question{
Establish the values of the parameter \(t\) 
which ensure that the equation 
\[ 
 x^{2} + tx + t + 8 = 0
\]
with an unknown \(x\) 
has complex solutions with a nonzero imaginary part.}
{\(\left (-4;8\right )\)}{\(\left [ -4;8\right ] \)}{\(\left (-\infty ;-4\right )\cup \left (8;\infty \right )\)}{\(\left (-\infty ;-4\right ] \cup \left [ 8;\infty \right )\)}{}{}}
\LANGpl{
\Question{
Określ wartości parametru  \(t\), które gwarantują, że równanie
\[ 
 x^{2} + tx + t + 8 = 0
\]
z niewiadomą  \(x\) 
ma złożone rozwiązania z umowną niezerową częścią.}
{\(\left (-4;8\right )\)}{\(\left [ -4;8\right ] \)}{\(\left (-\infty ;-4\right )\cup \left (8;\infty \right )\)}{\(\left (-\infty ;-4\right ] \cup \left [ 8;\infty \right )\)}{}{}}
\LANGcs{
\Question{
Množina všech takových parametrů
\(t\), pro než
má rovnice \(x^{2} + tx + t + 8 = 0\) 
s neznámou \(x\) 
imaginární kořeny (tj. komplexní kořeny s nenulovou imaginární částí), je:}
{\(\left (-4;8\right )\)}{\(\left \langle -4;8\right \rangle \)}{\(\left (-\infty ;-4\right )\cup \left (8;\infty \right )\)}{\(\left (-\infty ;-4\right \rangle \cup \left \langle 8;\infty \right )\)}{}{}}
\LANGsk{
\Question{
Množina všetkých takých parametrov \(t\), pre ktoré má rovnica 
 \[ 
 x^{2} + tx + t + 8 = 0
\]
s neznámou \(x\) 
imaginárne korene (tj. komplexné korene s nenulovú imaginárny časťou), je:}
{\(\left (-4;8\right )\)}{\(\left [ -4;8\right ] \)}{\(\left (-\infty ;-4\right )\cup \left (8;\infty \right )\)}{\(\left (-\infty ;-4\right ] \cup \left [ 8;\infty \right )\)}{}{}}
\LANGes{
\Question{
Determina los valores del parámetro \(t\) 
suponiendo que la ecuación
\[ 
 x^{2} + tx + t + 8 = 0
\]
con una incógnita \(x\) 
tiene soluciones complejas con una parte imaginaria distinta de cero.}
{\(\left (-4;8\right )\)}{\(\left [ -4;8\right ] \)}{\(\left (-\infty ;-4\right )\cup \left (8;\infty \right )\)}{\(\left (-\infty ;-4\right ] \cup \left [ 8;\infty \right )\)}{}{}}
\End

\Data\ID{9000035603}
\Updated{2020-07-18 02:07:37}
\Level{A}
\SubArea{Quadratic equations with complex roots}
\LANGen{
\Question{
Find the solution set of the following equation. 
\[ 
 4x^{2} + 9 = 0
\]}
{\(\left \{-\frac{3}
{2}\mathrm{i}; \frac{3} 
{2}\mathrm{i}\right \}\)}{\(\left \{-\frac{2}
{3}\mathrm{i}; \frac{2} 
{3}\mathrm{i}\right \}\)}{\(\left \{-\frac{9}
{4}\mathrm{i}; \frac{9} 
{4}\mathrm{i}\right \}\)}{\(\left \{-\frac{3}
{2}; \frac{3} 
{2}\right \}\)}{}{}}
\LANGpl{
\Question{
Wskaż zbiór, który jest rozwiązaniem następującego równania.
\[ 
 4x^{2} + 9 = 0
\]}
{\(\left \{-\frac{3}
{2}\mathrm{i}; \frac{3} 
{2}\mathrm{i}\right \}\)}{\(\left \{-\frac{2}
{3}\mathrm{i}; \frac{2} 
{3}\mathrm{i}\right \}\)}{\(\left \{-\frac{9}
{4}\mathrm{i}; \frac{9} 
{4}\mathrm{i}\right \}\)}{\(\left \{-\frac{3}
{2}; \frac{3} 
{2}\right \}\)}{}{}}
\LANGcs{
\Question{
Množina všech komplexních řešení rovnice
\(4x^{2} + 9 = 0\) je:}
{\(\left \{-\frac{3}
{2}\mathrm{i}; \frac{3} 
{2}\mathrm{i}\right \}\)}{\(\left \{-\frac{2}
{3}\mathrm{i}; \frac{2} 
{3}\mathrm{i}\right \}\)}{\(\left \{-\frac{9}
{4}\mathrm{i}; \frac{9} 
{4}\mathrm{i}\right \}\)}{\(\left \{-\frac{3}
{2}; \frac{3} 
{2}\right \}\)}{}{}}
\LANGsk{
\Question{
Nájdite množinu všetkých komplexných riešení danej rovnice. 
\[ 
 4x^{2} + 9 = 0 \]}
{\(\left \{-\frac{3}
{2}\mathrm{i}; \frac{3} 
{2}\mathrm{i}\right \}\)}{\(\left \{-\frac{2}
{3}\mathrm{i}; \frac{2} 
{3}\mathrm{i}\right \}\)}{\(\left \{-\frac{9}
{4}\mathrm{i}; \frac{9} 
{4}\mathrm{i}\right \}\)}{\(\left \{-\frac{3}
{2}; \frac{3} 
{2}\right \}\)}{}{}}
\LANGes{
\Question{
Determina el conjunto de soluciones de la siguiente ecuación.
\[ 
 4x^{2} + 9 = 0
\]}
{\(\left \{-\frac{3}
{2}\mathrm{i}; \frac{3} 
{2}\mathrm{i}\right \}\)}{\(\left \{-\frac{2}
{3}\mathrm{i}; \frac{2} 
{3}\mathrm{i}\right \}\)}{\(\left \{-\frac{9}
{4}\mathrm{i}; \frac{9} 
{4}\mathrm{i}\right \}\)}{\(\left \{-\frac{3}
{2}; \frac{3} 
{2}\right \}\)}{}{}}
\End

\Data\ID{9000039107}
\Updated{2020-07-18 02:07:21}
\Level{A}
\SubArea{Quadratic equations with complex roots}
\LANGen{
\Question{
Find the sum of the reciprocal values of the solutions of
\(5x^{2} - 2x + 1 = 0\).}
{\(2\)}{\(\frac{2}
{5}\)}{\(\frac{2}
{5} + \frac{4} 
{5}\mathrm{i}\)}{\(-\frac{2} 
{5}\)}{}{}}
\LANGpl{
\Question{
Wskaż sumę odwrotności rozwiązań 
\(5x^{2} - 2x + 1 = 0\).}
{\(2\)}{\(\frac{2}
{5}\)}{\(\frac{2}
{5} + \frac{4} 
{5}\mathrm{i}\)}{\(-\frac{2} 
{5}\)}{}{}}
\LANGcs{
\Question{
Součet převrácených hodnot kořenů kvadratické rovnice
\(5x^{2} - 2x + 1 = 0\) je:}
{\(2\)}{\(\frac{2}
{5}\)}{\(\frac{2}
{5} + \frac{4} 
{5}\mathrm{i}\)}{\(-\frac{2} 
{5}\)}{}{}}
\LANGsk{
\Question{
Súčet prevrátených hodnôt koreňov kvadratickej rovnice 
\(5x^{2} - 2x + 1 = 0\) je:}
{\(2\)}{\(\frac{2}
{5}\)}{\(\frac{2}
{5} + \frac{4} 
{5}\mathrm{i}\)}{\(-\frac{2} 
{5}\)}{}{}}
\LANGes{
\Question{
Determina la suma de los valores recíprocos de las soluciones de
\(5x^{2} - 2x + 1 = 0\).}
{\(2\)}{\(\frac{2}
{5}\)}{\(\frac{2}
{5} + \frac{4} 
{5}\mathrm{i}\)}{\(-\frac{2} 
{5}\)}{}{}}
\End

\Data\ID{9000064502}
\Updated{2020-07-18 02:07:09}
\Level{A}
\SubArea{Quadratic equations with complex roots}
\LANGen{
\Question{
Find the quadratic equation with the solution
\(x_{1, 2} = 2\pm \mathrm{i}\sqrt{2}\).}
{\(x^{2} - 4x + 6 = 0\)}{\(3x^{2} + 4x + 2 = 0\)}{\(3x^{2} - 4x + 2 = 0\)}{\(x^{2} + 4x + 6 = 0\)}{}{}}
\LANGpl{
\Question{
Wskaż równanie kwadratowe, którego rozwiązaniem jest 
\(x_{1, 2} = 2\pm \mathrm{i}\sqrt{2}\).}
{\(x^{2} - 4x + 6 = 0\)}{\(3x^{2} + 4x + 2 = 0\)}{\(3x^{2} - 4x + 2 = 0\)}{\(x^{2} + 4x + 6 = 0\)}{}{}}
\LANGcs{
\Question{
Komplexní čísla \(x_{1, 2} = 2\pm \mathrm{i}\sqrt{2}\) 
jsou kořeny kvadratické rovnice:}
{\(x^{2} - 4x + 6 = 0\)}{\(3x^{2} + 4x + 2 = 0\)}{\(3x^{2} - 4x + 2 = 0\)}{\(x^{2} + 4x + 6 = 0\)}{}{}}
\LANGsk{
\Question{
Nájdite kvadratickú rovnicu, ktorej korene sú komplexné čísla
\(x_{1, 2} = 2\pm \mathrm{i}\sqrt{2}\).}
{\(x^{2} - 4x + 6 = 0\)}{\(3x^{2} + 4x + 2 = 0\)}{\(3x^{2} - 4x + 2 = 0\)}{\(x^{2} + 4x + 6 = 0\)}{}{}}
\LANGes{
\Question{
Determina la ecuación cuadrática con la solución
\(x_{1, 2} = 2\pm \mathrm{i}\sqrt{2}\).}
{\(x^{2} - 4x + 6 = 0\)}{\(3x^{2} + 4x + 2 = 0\)}{\(3x^{2} - 4x + 2 = 0\)}{\(x^{2} + 4x + 6 = 0\)}{}{}}
\End

\Data\ID{9000064503}
\Updated{2020-07-18 02:07:17}
\Level{A}
\SubArea{Quadratic equations with complex roots}
\LANGen{
\Question{
Find the values of the real coefficients
\(a\),
\(b\) and
\(c\) such
that the quadratic equation 
\[ 
 ax^{2} + bx + c = 0
\]
has solution \(x_{1, 2} =\pm \mathrm{i}\frac{\sqrt{5}} 
 {3} \).}
{\(a = 9\text{, }b = 0\text{, }c = 5\)}{\(a = 5\text{, }b = 0\text{, }c = 9\)}{\(a = 9\text{, }b = 0\text{, }c = -5\)}{\(a = 5\text{, }b = 0\text{, }c = -9\)}{}{}}
\LANGpl{
\Question{
Wyznacz wartości współczynników rzeczywistych 
\(a\),
\(b\) i
\(c\) tak,
aby równanie kwadratowe
\[ 
 ax^{2} + bx + c = 0
\]
miało rozwiązania \(x_{1, 2} =\pm \mathrm{i}\frac{\sqrt{5}} 
 {3} \).}
{\(a = 9\text{, }b = 0\text{, }c = 5\)}{\(a = 5\text{, }b = 0\text{, }c = 9\)}{\(a = 9\text{, }b = 0\text{, }c = -5\)}{\(a = 5\text{, }b = 0\text{, }c = -9\)}{}{}}
\LANGcs{
\Question{
Kvadratická rovnice \(ax^{2} + bx + c = 0\) s
komplexními kořeny \(x_{1, 2} =\pm \mathrm{i}\frac{\sqrt{5}} 
 {3} \) 
má koeficienty:}
{\(a = 9\text{, }b = 0\text{, }c = 5\)}{\(a = 5\text{, }b = 0\text{, }c = 9\)}{\(a = 9\text{, }b = 0\text{, }c = -5\)}{\(a = 5\text{, }b = 0\text{, }c = -9\)}{}{}}
\LANGsk{
\Question{
Nájdite hodnoty reálnych koeficientov
\(a\),
\(b\) a
\(c\) tak, aby
kvadratická rovnica  
\[ 
 ax^{2} + bx + c = 0
\]
mala komplexné korene \(x_{1, 2} =\pm \mathrm{i}\frac{\sqrt{5}} 
 {3} \).}
{\(a = 9\text{, }b = 0\text{, }c = 5\)}{\(a = 5\text{, }b = 0\text{, }c = 9\)}{\(a = 9\text{, }b = 0\text{, }c = -5\)}{\(a = 5\text{, }b = 0\text{, }c = -9\)}{}{}}
\LANGes{
\Question{
Determina los valores de coeficientes reales
\(a\),
\(b\) y
\(c\) suponiendo que la ecuación cuadrática 
\[ 
 ax^{2} + bx + c = 0
\]
tiene la solución \(x_{1, 2} =\pm \mathrm{i}\frac{\sqrt{5}} 
 {3} \).}
{\(a = 9\text{, }b = 0\text{, }c = 5\)}{\(a = 5\text{, }b = 0\text{, }c = 9\)}{\(a = 9\text{, }b = 0\text{, }c = -5\)}{\(a = 5\text{, }b = 0\text{, }c = -9\)}{}{}}
\End

\Data\ID{9000064504}
\Updated{2020-07-18 02:07:09}
\Level{A}
\SubArea{Quadratic equations with complex roots}
\LANGen{
\Question{
Find the values of the real coefficients
\(a\),
\(b\) and
\(c\) such
that the quadratic equation 
\[ 
 ax^{2} + bx + c = 0
\]
has solutions \(x_{1, 2} = 1\pm \frac{\mathrm{i}}
{2}\).}
{\(a = 4\text{, }b = -8\text{, }c = 5\)}{\(a = 1\text{, }b = -4\text{, }c = 5\)}{\(a = 4\text{, }b = 8\text{, }c = 5\)}{\(a = 1\text{, }b = 4\text{, }c = 5\)}{}{}}
\LANGpl{
\Question{
Wskaż wartości współczynników rzeczywistych
\(a\),
\(b\) i
\(c\) tak, aby
równanie kwadratowe
\[ 
 ax^{2} + bx + c = 0
\]
miało rozwiązania  \(x_{1, 2} = 1\pm \frac{\mathrm{i}}
{2}\).}
{\(a = 4\text{, }b = -8\text{, }c = 5\)}{\(a = 1\text{, }b = -4\text{, }c = 5\)}{\(a = 4\text{, }b = 8\text{, }c = 5\)}{\(a = 1\text{, }b = 4\text{, }c = 5\)}{}{}}
\LANGcs{
\Question{
Kvadratická rovnice \(ax^{2} + bx + c = 0\) s
komplexními kořeny \(x_{1, 2} = 1\pm \frac{\mathrm{i}}
{2}\) 
má koeficienty:}
{\(a = 4\text{, }b = -8\text{, }c = 5\)}{\(a = 1\text{, }b = -4\text{, }c = 5\)}{\(a = 4\text{, }b = 8\text{, }c = 5\)}{\(a = 1\text{, }b = 4\text{, }c = 5\)}{}{}}
\LANGsk{
\Question{
Nájdite hodnoty reálnych koeficientov
\(a\),
\(b\) a
\(c\) tak, aby
kvadratická rovnica 
\[ 
 ax^{2} + bx + c = 0
\]
mala komplexné korene \(x_{1, 2} = 1\pm \frac{\mathrm{i}}
{2}\).}
{\(a = 4\text{, }b = -8\text{, }c = 5\)}{\(a = 1\text{, }b = -4\text{, }c = 5\)}{\(a = 4\text{, }b = 8\text{, }c = 5\)}{\(a = 1\text{, }b = 4\text{, }c = 5\)}{}{}}
\LANGes{
\Question{
Determina los valores de los coeficientes reales
\(a\),
\(b\) y
\(c\) suponiendo que la ecuación cuadrática
\[ 
 ax^{2} + bx + c = 0
\]
tiene soluciones \(x_{1, 2} = 1\pm \frac{\mathrm{i}}
{2}\).}
{\(a = 4\text{, }b = -8\text{, }c = 5\)}{\(a = 1\text{, }b = -4\text{, }c = 5\)}{\(a = 4\text{, }b = 8\text{, }c = 5\)}{\(a = 1\text{, }b = 4\text{, }c = 5\)}{}{}}
\End

\Data\ID{9000064505}
\Updated{2020-07-18 02:07:18}
\Level{A}
\SubArea{Quadratic equations with complex roots}
\LANGen{
\Question{
Find the factorization of the following quadratic polynomial in the set of polynomial
with complex valued coefficients. 
\[ 
 2x^{2} + 32
\]}
{\(2(x + 4\mathrm{i})(x - 4\mathrm{i})\)}{\(2(x - 4\mathrm{i})^{2}\)}{\((x + 4\mathrm{i})(x - 4\mathrm{i})\)}{\(2(x + 4\mathrm{i})^{2}\)}{}{}}
\LANGpl{
\Question{
Rozłóż wielomian kwadratowy na czynniki pierwsze.
\[ 
 2x^{2} + 32
\]}
{\(2(x + 4\mathrm{i})(x - 4\mathrm{i})\)}{\(2(x - 4\mathrm{i})^{2}\)}{\((x + 4\mathrm{i})(x - 4\mathrm{i})\)}{\(2(x + 4\mathrm{i})^{2}\)}{}{}}
\LANGcs{
\Question{
Kvadratický výraz \(2x^{2} + 32\) 
lze v množině komplexních čísel rozložit na součin:}
{\(2(x + 4\mathrm{i})(x - 4\mathrm{i})\)}{\(2(x - 4\mathrm{i})^{2}\)}{\((x + 4\mathrm{i})(x - 4\mathrm{i})\)}{\(2(x + 4\mathrm{i})^{2}\)}{}{}}
\LANGsk{
\Question{
Kvadratický výraz 
\[ 
 2x^{2} + 32
\] môžeme v množine komplexných čísel rozložiť na súčin koreňových činiteľov:}
{\(2(x + 4\mathrm{i})(x - 4\mathrm{i})\)}{\(2(x - 4\mathrm{i})^{2}\)}{\((x + 4\mathrm{i})(x - 4\mathrm{i})\)}{\(2(x + 4\mathrm{i})^{2}\)}{}{}}
\LANGes{
\Question{
Determina la factorización del siguiente polinomio cuadrático en el conjunto de polinomios.
con coeficientes  complejos.
\[ 
 2x^{2} + 32
\]}
{\(2(x + 4\mathrm{i})(x - 4\mathrm{i})\)}{\(2(x - 4\mathrm{i})^{2}\)}{\((x + 4\mathrm{i})(x - 4\mathrm{i})\)}{\(2(x + 4\mathrm{i})^{2}\)}{}{}}
\End

\Data\ID{9000064506}
\Updated{2020-07-18 09:07:04}
\Level{A}
\SubArea{Quadratic equations with complex roots}
\LANGen{
\Question{
Find the factorization of the following quadratic polynomial in the set of polynomial
with complex valued coefficients. 
\[ 
 2x^{2} + 4x + 5
\]}
{\(2\! \left (x + 1 + \frac{\sqrt{6}} 
 {2} \mathrm{i}\right )\! \! \left (x + 1 -\frac{\sqrt{6}} 
 {2} \mathrm{i}\right )\)}{\(2\! \left (x - 1 + \frac{\sqrt{6}} 
 {2} \mathrm{i}\right )\! \! \left (x - 1 -\frac{\sqrt{6}} 
 {2} \mathrm{i}\right )\)}{\(\left (x + 1 -\frac{\sqrt{6}} 
 {2} \mathrm{i}\right )\! \! \left (x + 1 + \frac{\sqrt{6}} 
 {2} \mathrm{i}\right )\)}{\(\left (x - 1 -\frac{\sqrt{6}} 
 {2} \mathrm{i}\right )\! \! \left (x - 1 + \frac{\sqrt{6}} 
 {2} \mathrm{i}\right )\)}{}{}}
\LANGpl{
\Question{
Rozłóż wielomian kwadratowy na czynniki pierwsze.
\[ 
 2x^{2} + 4x + 5
\]}
{\(2\! \left (x + 1 + \frac{\sqrt{6}} 
 {2} \mathrm{i}\right )\! \! \left (x + 1 -\frac{\sqrt{6}} 
 {2} \mathrm{i}\right )\)}{\(2\! \left (x - 1 + \frac{\sqrt{6}} 
 {2} \mathrm{i}\right )\! \! \left (x - 1 -\frac{\sqrt{6}} 
 {2} \mathrm{i}\right )\)}{\(\left (x + 1 -\frac{\sqrt{6}} 
 {2} \mathrm{i}\right )\! \! \left (x + 1 + \frac{\sqrt{6}} 
 {2} \mathrm{i}\right )\)}{\(\left (x - 1 -\frac{\sqrt{6}} 
 {2} \mathrm{i}\right )\! \! \left (x - 1 + \frac{\sqrt{6}} 
 {2} \mathrm{i}\right )\)}{}{}}
\LANGcs{
\Question{
Kvadratický trojčlen \[2x^{2} + 4x + 5\] můžeme v množině komplexních čísel rozložit na součin kořenových činitelů:}
{\(2\! \left (x + 1 + \frac{\sqrt{6}} 
 {2} \mathrm{i}\right )\! \! \left (x + 1 -\frac{\sqrt{6}} 
 {2} \mathrm{i}\right )\)}{\(2\! \left (x - 1 + \frac{\sqrt{6}} 
 {2} \mathrm{i}\right )\! \! \left (x - 1 -\frac{\sqrt{6}} 
 {2} \mathrm{i}\right )\)}{\(\left (x + 1 -\frac{\sqrt{6}} 
 {2} \mathrm{i}\right )\! \! \left (x + 1 + \frac{\sqrt{6}} 
 {2} \mathrm{i}\right )\)}{\(\left (x - 1 -\frac{\sqrt{6}} 
 {2} \mathrm{i}\right )\! \! \left (x - 1 + \frac{\sqrt{6}} 
 {2} \mathrm{i}\right )\)}{}{}}
\LANGsk{
\Question{
Kvadratický trojčlen 
\[ 
 2x^{2} + 4x + 5
\] môžeme v množine komplexných čísel rozložiť na súčin koreňových činiteľov:}
{\(2\! \left (x + 1 + \frac{\sqrt{6}} 
 {2} \mathrm{i}\right )\! \! \left (x + 1 -\frac{\sqrt{6}} 
 {2} \mathrm{i}\right )\)}{\(2\! \left (x - 1 + \frac{\sqrt{6}} 
 {2} \mathrm{i}\right )\! \! \left (x - 1 -\frac{\sqrt{6}} 
 {2} \mathrm{i}\right )\)}{\(\left (x + 1 -\frac{\sqrt{6}} 
 {2} \mathrm{i}\right )\! \! \left (x + 1 + \frac{\sqrt{6}} 
 {2} \mathrm{i}\right )\)}{\(\left (x - 1 -\frac{\sqrt{6}} 
 {2} \mathrm{i}\right )\! \! \left (x - 1 + \frac{\sqrt{6}} 
 {2} \mathrm{i}\right )\)}{}{}}
\LANGes{
\Question{
Determina la factorización del siguiente polinomio cuadrático en el conjunto de polinomios. con coeficientes complejos.
\[ 
 2x^{2} + 4x + 5
\]}
{\(2\! \left (x + 1 + \frac{\sqrt{6}} 
 {2} \mathrm{i}\right )\! \! \left (x + 1 -\frac{\sqrt{6}} 
 {2} \mathrm{i}\right )\)}{\(2\! \left (x - 1 + \frac{\sqrt{6}} 
 {2} \mathrm{i}\right )\! \! \left (x - 1 -\frac{\sqrt{6}} 
 {2} \mathrm{i}\right )\)}{\(\left (x + 1 -\frac{\sqrt{6}} 
 {2} \mathrm{i}\right )\! \! \left (x + 1 + \frac{\sqrt{6}} 
 {2} \mathrm{i}\right )\)}{\(\left (x - 1 -\frac{\sqrt{6}} 
 {2} \mathrm{i}\right )\! \! \left (x - 1 + \frac{\sqrt{6}} 
 {2} \mathrm{i}\right )\)}{}{}}
\End

\Data\ID{9000064507}
\Updated{2020-07-18 09:07:23}
\Level{A}
\SubArea{Quadratic equations with complex roots}
\LANGen{
\Question{
Solve the following quadratic equation in the complex plane. 
\[ 
 4x^{2} + 12 = 0
\]}
{\(x_{1, 2} =\pm \mathrm{i}\sqrt{3}\)}{\(x_{1, 2} =\pm 3\)}{\(x_{1, 2} =\pm 3\mathrm{i}\)}{\(x_{1, 2} =\pm \sqrt{3}\)}{}{}}
\LANGpl{
\Question{
Rozwiąż równanie kwadratowe na płaszczyźnie zespolonej.
\[ 
 4x^{2} + 12 = 0
\]}
{\(x_{1, 2} =\pm \mathrm{i}\sqrt{3}\)}{\(x_{1, 2} =\pm 3\)}{\(x_{1, 2} =\pm 3\mathrm{i}\)}{\(x_{1, 2} =\pm \sqrt{3}\)}{}{}}
\LANGcs{
\Question{
Vyřešte danou kvadratickou rovnici v množině komplexních čísel. \[4x^{2} + 12 = 0\]}
{\(x_{1, 2} =\pm \mathrm{i}\sqrt{3}\)}{\(x_{1, 2} =\pm 3\)}{\(x_{1, 2} =\pm 3\mathrm{i}\)}{\(x_{1, 2} =\pm \sqrt{3}\)}{}{}}
\LANGsk{
\Question{
Vyriešte danú kvadratickú rovnicu v množine komplexných čísel. 
\[ 
 4x^{2} + 12 = 0
\]}
{\(x_{1, 2} =\pm \mathrm{i}\sqrt{3}\)}{\(x_{1, 2} =\pm 3\)}{\(x_{1, 2} =\pm 3\mathrm{i}\)}{\(x_{1, 2} =\pm \sqrt{3}\)}{}{}}
\LANGes{
\Question{
Resuelve la siguiente ecuación cuadrática en el plano complejo.
\[ 
 4x^{2} + 12 = 0
\]}
{\(x_{1, 2} =\pm \mathrm{i}\sqrt{3}\)}{\(x_{1, 2} =\pm 3\)}{\(x_{1, 2} =\pm 3\mathrm{i}\)}{\(x_{1, 2} =\pm \sqrt{3}\)}{}{}}
\End

\Data\ID{9000064508}
\Updated{2020-07-18 09:07:51}
\Level{A}
\SubArea{Quadratic equations with complex roots}
\LANGen{
\Question{
Solve the following quadratic equation in the complex plane. 
\[ 
 2x^{2} + x + 1 = 0
\]}
{\(x_{1, 2} = \frac{-1\pm \mathrm{i}\sqrt{7}} 
 {4} \)}{\(x_{1, 2} = \frac{-1\pm \mathrm{i}\sqrt{7}} 
 {2} \)}{\(x_{1, 2} = \frac{1\pm \mathrm{i}\sqrt{7}} 
 {4} \)}{\(x_{1, 2} = \frac{1\pm \mathrm{i}\sqrt{7}} 
 {2} \)}{}{}}
\LANGpl{
\Question{
Rozwiąż równanie kwadratowe na płaszczyźnie zespolonej.
\[ 
 2x^{2} + x + 1 = 0
\]}
{\(x_{1, 2} = \frac{-1\pm \mathrm{i}\sqrt{7}} 
 {4} \)}{\(x_{1, 2} = \frac{-1\pm \mathrm{i}\sqrt{7}} 
 {2} \)}{\(x_{1, 2} = \frac{1\pm \mathrm{i}\sqrt{7}} 
 {4} \)}{\(x_{1, 2} = \frac{1\pm \mathrm{i}\sqrt{7}} 
 {2} \)}{}{}}
\LANGcs{
\Question{
Vyřešte danou kvadratickou rovnici v množině komplexních čísel. \[2x^{2} + x + 1 = 0\]}
{\(x_{1, 2} = \frac{-1\pm \mathrm{i}\sqrt{7}} 
 {4} \)}{\(x_{1, 2} = \frac{-1\pm \mathrm{i}\sqrt{7}} 
 {2} \)}{\(x_{1, 2} = \frac{1\pm \mathrm{i}\sqrt{7}} 
 {4} \)}{\(x_{1, 2} = \frac{1\pm \mathrm{i}\sqrt{7}} 
 {2} \)}{}{}}
\LANGsk{
\Question{
Vyriešte danú kvadratickú rovnicu v množine komplexných čísel. 
\[ 
 2x^{2} + x + 1 = 0
\]}
{\(x_{1, 2} = \frac{-1\pm \mathrm{i}\sqrt{7}} 
 {4} \)}{\(x_{1, 2} = \frac{-1\pm \mathrm{i}\sqrt{7}} 
 {2} \)}{\(x_{1, 2} = \frac{1\pm \mathrm{i}\sqrt{7}} 
 {4} \)}{\(x_{1, 2} = \frac{1\pm \mathrm{i}\sqrt{7}} 
 {2} \)}{}{}}
\LANGes{
\Question{
Resuelve la siguiente ecuación cuadrática en el plano complejo.
\[ 
 2x^{2} + x + 1 = 0
\]}
{\(x_{1, 2} = \frac{-1\pm \mathrm{i}\sqrt{7}} 
 {4} \)}{\(x_{1, 2} = \frac{-1\pm \mathrm{i}\sqrt{7}} 
 {2} \)}{\(x_{1, 2} = \frac{1\pm \mathrm{i}\sqrt{7}} 
 {4} \)}{\(x_{1, 2} = \frac{1\pm \mathrm{i}\sqrt{7}} 
 {2} \)}{}{}}
\End

\Data\ID{9000069901}
\Updated{2020-07-18 09:07:58}
\Level{A}
\SubArea{Quadratic equations with complex roots}
\LANGen{
\Question{
Solve the following quadratic equation in the complex plane. 
\[ 
 x^{2} + 4x + 5 = 0
\]}
{\(x_{1} = -2 + \mathrm{i}\),
\(x_{2} = -2 -\mathrm{i}\)}{\(x = -2\)}{\(x_{1} = 2 + \mathrm{i}\),
\(x_{2} = 2 -\mathrm{i}\)}{\(x_{1} = -3\),
\(x_{2} = -1\)}{}{}}
\LANGpl{
\Question{
Rozwiąż równanie kwadratowe na płaszczyźnie zespolonej.
\[ 
 x^{2} + 4x + 5 = 0
\]}
{\(x_{1} = -2 + \mathrm{i}\),
\(x_{2} = -2 -\mathrm{i}\)}{\(x = -2\)}{\(x_{1} = 2 + \mathrm{i}\),
\(x_{2} = 2 -\mathrm{i}\)}{\(x_{1} = -3\),
\(x_{2} = -1\)}{}{}}
\LANGcs{
\Question{
Kvadratická rovnice \(x^{2} + 4x + 5 = 0\) 
řešená v množině komplexních čísel má kořeny:}
{\(x_{1} = -2 + \mathrm{i}\), \( x_{2} = -2 -\mathrm{i}\)}{\(x = -2\)}{\(x_{1} = 2 + \mathrm{i}\), \( x_{2} = 2 -\mathrm{i}\)}{\(x_{1} = -3\), \( x_{2} = -1\)}{}{}}
\LANGsk{
\Question{
Vyriešte danú kvadratickú rovnicu v množine komplexných čísel. 
\[ 
 x^{2} + 4x + 5 = 0
\]}
{\(x_{1} = -2 + \mathrm{i}\),
\(x_{2} = -2 -\mathrm{i}\)}{\(x = -2\)}{\(x_{1} = 2 + \mathrm{i}\),
\(x_{2} = 2 -\mathrm{i}\)}{\(x_{1} = -3\),
\(x_{2} = -1\)}{}{}}
\LANGes{
\Question{
Resuelve la siguiente ecuación cuadrática en el plano complejo.
\[ 
 x^{2} + 4x + 5 = 0
\]}
{\(x_{1} = -2 + \mathrm{i}\),
\(x_{2} = -2 -\mathrm{i}\)}{\(x = -2\)}{\(x_{1} = 2 + \mathrm{i}\),
\(x_{2} = 2 -\mathrm{i}\)}{\(x_{1} = -3\),
\(x_{2} = -1\)}{}{}}
\End

\Data\ID{9000069902}
\Updated{2020-07-18 09:07:13}
\Level{A}
\SubArea{Quadratic equations with complex roots}
\LANGen{
\Question{
Solve the following quadratic equation in the complex plane. 
\[ 
 3x^{2} + 2x + 2 = 0
\]}
{\(x_{1} = -\frac{1}
{3} + \frac{\sqrt{5}} 
 {3} \mathrm{i}\),
\(x_{2} = -\frac{1}
{3} -\frac{\sqrt{5}} 
 {3} \mathrm{i}\)}{\(x_{1} = -\frac{1}
{3}\)}{\(x_{1} = \frac{1} 
{3} + \frac{\sqrt{5}} 
 {3} \),
\(x_{2} = \frac{1} 
{3} + \frac{\sqrt{5}} 
 {3} \)}{\(x_{1} = \frac{1} 
{3} + \frac{\sqrt{5}} 
 {3} \mathrm{i}\),
\(x_{2} = \frac{1} 
{3} -\frac{\sqrt{5}} 
 {3} \mathrm{i}\)}{}{}}
\LANGpl{
\Question{
Rozwiąż równanie kwadratowe na płaszczyźnie zespolonej. 
\[ 
 3x^{2} + 2x + 2 = 0
\]}
{\(x_{1} = -\frac{1}
{3} + \frac{\sqrt{5}} 
 {3} \mathrm{i}\),
\(x_{2} = -\frac{1}
{3} -\frac{\sqrt{5}} 
 {3} \mathrm{i}\)}{\(x_{1} = -\frac{1}
{3}\)}{\(x_{1} = \frac{1} 
{3} + \frac{\sqrt{5}} 
 {3} \),
\(x_{2} = \frac{1} 
{3} + \frac{\sqrt{5}} 
 {3} \)}{\(x_{1} = \frac{1} 
{3} + \frac{\sqrt{5}} 
 {3} \mathrm{i}\),
\(x_{2} = \frac{1} 
{3} -\frac{\sqrt{5}} 
 {3} \mathrm{i}\)}{}{}}
\LANGcs{
\Question{
Kvadratická rovnice \(3x^{2} + 2x + 2 = 0\) 
řešená v množině komplexních čísel má kořeny:}
{\(x_{1} = -\frac{1}
{3} + \frac{\sqrt{5}} 
 {3} \mathrm{i}\), \( x_{2} = -\frac{1}
{3} -\frac{\sqrt{5}} 
 {3} \mathrm{i}\)}{\(x_{1} = -\frac{1}
{3}\)}{\(x_{1} = \frac{1} 
{3} + \frac{\sqrt{5}} 
 {3} \), \( x_{2} = \frac{1} 
{3} + \frac{\sqrt{5}} 
 {3} \)}{\(x_{1} = \frac{1} 
{3} + \frac{\sqrt{5}} 
 {3} \mathrm{i}\), \(x_{2} = \frac{1} 
{3} -\frac{\sqrt{5}} 
 {3} \mathrm{i}\)}{}{}}
\LANGsk{
\Question{
Vyriešte danú kvadratickú rovnicu v množine komplexných čísel. 
\[ 
 3x^{2} + 2x + 2 = 0
\]}
{\(x_{1} = -\frac{1}
{3} + \frac{\sqrt{5}} 
 {3} \mathrm{i}\),
\(x_{2} = -\frac{1}
{3} -\frac{\sqrt{5}} 
 {3} \mathrm{i}\)}{\(x_{1} = -\frac{1}
{3}\)}{\(x_{1} = \frac{1} 
{3} + \frac{\sqrt{5}} 
 {3} \),
\(x_{2} = \frac{1} 
{3} + \frac{\sqrt{5}} 
 {3} \)}{\(x_{1} = \frac{1} 
{3} + \frac{\sqrt{5}} 
 {3} \mathrm{i}\),
\(x_{2} = \frac{1} 
{3} -\frac{\sqrt{5}} 
 {3} \mathrm{i}\)}{}{}}
\LANGes{
\Question{
Resuelve la siguiente ecuación cuadrática en el plano complejo.
\[ 
 3x^{2} + 2x + 2 = 0
\]}
{\(x_{1} = -\frac{1}
{3} + \frac{\sqrt{5}} 
 {3} \mathrm{i}\),
\(x_{2} = -\frac{1}
{3} -\frac{\sqrt{5}} 
 {3} \mathrm{i}\)}{\(x_{1} = -\frac{1}
{3}\)}{\(x_{1} = \frac{1} 
{3} + \frac{\sqrt{5}} 
 {3} \),
\(x_{2} = \frac{1} 
{3} + \frac{\sqrt{5}} 
 {3} \)}{\(x_{1} = \frac{1} 
{3} + \frac{\sqrt{5}} 
 {3} \mathrm{i}\),
\(x_{2} = \frac{1} 
{3} -\frac{\sqrt{5}} 
 {3} \mathrm{i}\)}{}{}}
\End

\Data\ID{9000069903}
\Updated{2020-07-18 09:07:55}
\Level{A}
\SubArea{Quadratic equations with complex roots}
\LANGen{
\Question{
Find the factorization of the quadratic polynomial 
\[ 
 x^{2} + 2x + 2
\]
in the set of polynomials with complex valued coefficients.}
{\((x + 1 + \mathrm{i})(x + 1 -\mathrm{i})\)}{\((x - 1 + \mathrm{i})(x - 1 -\mathrm{i})\)}{\((x -\mathrm{i})(x + \mathrm{i})\)}{\((x - 1 + \mathrm{i})(x + 1 -\mathrm{i})\)}{}{}}
\LANGpl{
\Question{
Rozłóż na czynniki pierwsze wielomian kwadratowy
\[ 
 x^{2} + 2x + 2
\]
w zbiorze wielomianów o zespolonych współczynnikach rzeczywistych.}
{\((x + 1 + \mathrm{i})(x + 1 -\mathrm{i})\)}{\((x - 1 + \mathrm{i})(x - 1 -\mathrm{i})\)}{\((x -\mathrm{i})(x + \mathrm{i})\)}{\((x - 1 + \mathrm{i})(x + 1 -\mathrm{i})\)}{}{}}
\LANGcs{
\Question{
Kvadratický trojčlen \(x^{2} + 2x + 2\) 
můžeme v množině \(\mathbb{C}\) 
rozložit na součin kořenových činitelů:}
{\((x + 1 + \mathrm{i})(x + 1 -\mathrm{i})\)}{\((x - 1 + \mathrm{i})(x - 1 -\mathrm{i})\)}{\((x -\mathrm{i})(x + \mathrm{i})\)}{\((x - 1 + \mathrm{i})(x + 1 -\mathrm{i})\)}{}{}}
\LANGsk{
\Question{
Kvadratický trojčlen 
\[ 
 x^{2} + 2x + 2
\]
môžeme v množine komplexných čísel rozložiť na súčin koreňových činiteľov:}
{\((x + 1 + \mathrm{i})(x + 1 -\mathrm{i})\)}{\((x - 1 + \mathrm{i})(x - 1 -\mathrm{i})\)}{\((x -\mathrm{i})(x + \mathrm{i})\)}{\((x - 1 + \mathrm{i})(x + 1 -\mathrm{i})\)}{}{}}
\LANGes{
\Question{
Determina la factorización del polinomio cuadrático
\[ 
 x^{2} + 2x + 2
\]
en el conjunto de polinomios con coeficientes complejos.}
{\((x + 1 + \mathrm{i})(x + 1 -\mathrm{i})\)}{\((x - 1 + \mathrm{i})(x - 1 -\mathrm{i})\)}{\((x -\mathrm{i})(x + \mathrm{i})\)}{\((x - 1 + \mathrm{i})(x + 1 -\mathrm{i})\)}{}{}}
\End

\Data\ID{9000069904}
\Updated{2020-07-18 09:07:58}
\Level{A}
\SubArea{Quadratic equations with complex roots}
\LANGen{
\Question{
Find the factorization of the quadratic polynomial 
\[ 
 x^{2} + 2x + 5
\]
in the set of polynomials with complex valued coefficients.}
{\((x + 1 - 2\mathrm{i})(x + 1 + 2\mathrm{i})\)}{\((x - 1 - 2\mathrm{i})(x - 1 + 2\mathrm{i})\)}{\((x + 1 - 2\mathrm{i})(x - 1 + 2\mathrm{i})\)}{\((x - 1 - 2\mathrm{i})(x + 1 + 2\mathrm{i})\)}{}{}}
\LANGpl{
\Question{
Rozłóż na czynniki pierwsze wielomian kwadratowy
\[ 
 x^{2} + 2x + 5
\]
w zbiorze wielomianów o zespolonych współczynnikach rzeczywistych.}
{\((x + 1 - 2\mathrm{i})(x + 1 + 2\mathrm{i})\)}{\((x - 1 - 2\mathrm{i})(x - 1 + 2\mathrm{i})\)}{\((x + 1 - 2\mathrm{i})(x - 1 + 2\mathrm{i})\)}{\((x - 1 - 2\mathrm{i})(x + 1 + 2\mathrm{i})\)}{}{}}
\LANGcs{
\Question{
Kvadratický trojčlen \(x^{2} + 2x + 5\) 
můžeme v množině \(\mathbb{C}\) 
rozložit na součin kořenových činitelů:}
{\((x + 1 - 2\mathrm{i})(x + 1 + 2\mathrm{i})\)}{\((x - 1 - 2\mathrm{i})(x - 1 + 2\mathrm{i})\)}{\((x + 1 - 2\mathrm{i})(x - 1 + 2\mathrm{i})\)}{\((x - 1 - 2\mathrm{i})(x + 1 + 2\mathrm{i})\)}{}{}}
\LANGsk{
\Question{
Kvadratický trojčlen  
\[ 
 x^{2} + 2x + 5
\]
môžeme v množine komplexných čísel rozložiť na súčin koreňových činiteľov:}
{\((x + 1 - 2\mathrm{i})(x + 1 + 2\mathrm{i})\)}{\((x - 1 - 2\mathrm{i})(x - 1 + 2\mathrm{i})\)}{\((x + 1 - 2\mathrm{i})(x - 1 + 2\mathrm{i})\)}{\((x - 1 - 2\mathrm{i})(x + 1 + 2\mathrm{i})\)}{}{}}
\LANGes{
\Question{
Determina la factorización del polinomio cuadrático
\[ 
 x^{2} + 2x + 5
\]
en el conjunto de polinomios con coeficientes complejos.}
{\((x + 1 - 2\mathrm{i})(x + 1 + 2\mathrm{i})\)}{\((x - 1 - 2\mathrm{i})(x - 1 + 2\mathrm{i})\)}{\((x + 1 - 2\mathrm{i})(x - 1 + 2\mathrm{i})\)}{\((x - 1 - 2\mathrm{i})(x + 1 + 2\mathrm{i})\)}{}{}}
\End

\Data\ID{9000069905}
\Updated{2020-07-18 09:07:12}
\Level{A}
\SubArea{Quadratic equations with complex roots}
\LANGen{
\Question{
Find the sum of all the complex solutions of the following quadratic equation.
\[ 
 5x^{2} + 4x + 8 = 0
\]}
{\(-\frac{4} 
{5}\)}{\(- \frac{4} 
{10}\)}{\(\frac{24}
 {5} \mathrm{i}\)}{\(0\)}{}{}}
\LANGpl{
\Question{
Wskaż sumę wszystkich rozwiązań zespolonych równania kwadratowego.
\[ 
 5x^{2} + 4x + 8 = 0
\]}
{\(-\frac{4} 
{5}\)}{\(- \frac{4} 
{10}\)}{\(\frac{24}
 {5} \mathrm{i}\)}{\(0\)}{}{}}
\LANGcs{
\Question{
Součet kořenů kvadratické rovnice
\(5x^{2} + 4x + 8 = 0\) 
řešené v oboru komplexních čísel je roven:}
{\(-\frac{4} 
{5}\)}{\(- \frac{4} 
{10}\)}{\(\frac{24}
 {5} \mathrm{i}\)}{\(0\)}{}{}}
\LANGsk{
\Question{
Súčet všetkých koreňov kvadratickej rovnice
\[ 
 5x^{2} + 4x + 8 = 0
\] v množine komplexných čísel sa rovná:}
{\(-\frac{4} 
{5}\)}{\(- \frac{4} 
{10}\)}{\(\frac{24}
 {5} \mathrm{i}\)}{\(0\)}{}{}}
\LANGes{
\Question{
Calcula la suma de todas las soluciones complejas de la siguiente ecuación cuadrática.
\[ 
 5x^{2} + 4x + 8 = 0
\]}
{\(-\frac{4} 
{5}\)}{\(- \frac{4} 
{10}\)}{\(\frac{24}
 {5} \mathrm{i}\)}{\(0\)}{}{}}
\End

\Data\ID{9000069906}
\Updated{2020-07-18 09:07:17}
\Level{A}
\SubArea{Quadratic equations with complex roots}
\LANGen{
\Question{
Find the sum of all the complex solutions of the following quadratic equation.
\[ 
 x^{2} - 8x + 17 = 0
\]}
{\(8\)}{\(4\)}{\(4\mathrm{i}\)}{\(0\)}{}{}}
\LANGpl{
\Question{
Wskaż sumę wszystkich rozwiązań zespolonych równania kwadratowego.
\[ 
 x^{2} - 8x + 17 = 0
\]}
{\(8\)}{\(4\)}{\(4\mathrm{i}\)}{\(0\)}{}{}}
\LANGcs{
\Question{
Součet kořenů kvadratické rovnice
\[x^{2} - 8x + 17 = 0\] 
řešené v oboru komplexních čísel je roven:}
{\(8\)}{\(4\)}{\(4\mathrm{i}\)}{\(0\)}{}{}}
\LANGsk{
\Question{
Súčet všetkých koreňov kvadratickej rovnice
\[ 
 x^{2} - 8x + 17 = 0
\]  v množine komplexných čísel sa rovná:}
{\(8\)}{\(4\)}{\(4\mathrm{i}\)}{\(0\)}{}{}}
\LANGes{
\Question{
Calcula la suma de todas las soluciones complejas de la siguiente ecuación cuadrática.
\[ 
 x^{2} - 8x + 17 = 0
\]}
{\(8\)}{\(4\)}{\(4\mathrm{i}\)}{\(0\)}{}{}}
\End

\Data\ID{9000069907}
\Updated{2020-07-18 09:07:10}
\Level{A}
\SubArea{Quadratic equations with complex roots}
\LANGen{
\Question{
Find the quadratic equation with real valued coefficients and one of the solutions
\(x_{1} = -5 + \mathrm{i}\).}
{\(x^{2} + 10x + 26 = 0\)}{\(x^{2} - 10x + 26 = 0\)}{\(x^{2} - 10x - 24 = 0\)}{\(x^{2} + 10x + 24 = 0\)}{}{}}
\LANGpl{
\Question{
Wskaż równanie kwadratowe o współczynnikach rzeczywistych, którego jednym z rozwiązań jest
\(x_{1} = -5 + \mathrm{i}\).}
{\(x^{2} + 10x + 26 = 0\)}{\(x^{2} - 10x + 26 = 0\)}{\(x^{2} - 10x - 24 = 0\)}{\(x^{2} + 10x + 24 = 0\)}{}{}}
\LANGcs{
\Question{
Kvadratickou rovnici s reálnými koeficienty, jejíž jeden kořen je
\(x_1=- 5 + \mathrm{i}\),
můžeme zapsat ve tvaru:}
{\(x^{2} + 10x + 26 = 0\)}{\(x^{2} - 10x + 26 = 0\)}{\(x^{2} - 10x - 24 = 0\)}{\(x^{2} + 10x + 24 = 0\)}{}{}}
\LANGsk{
\Question{
Nájdite kvadratickú rovnicu s reálnymi koeficientmi a jedným koreňom
\(x_{1} = -5 + \mathrm{i}\).}
{\(x^{2} + 10x + 26 = 0\)}{\(x^{2} - 10x + 26 = 0\)}{\(x^{2} - 10x - 24 = 0\)}{\(x^{2} + 10x + 24 = 0\)}{}{}}
\LANGes{
\Question{
Determina la ecuación cuadrática con coeficientes reales suponiendo que una de las soluciones es:
\(x_{1} = -5 + \mathrm{i}\).}
{\(x^{2} + 10x + 26 = 0\)}{\(x^{2} - 10x + 26 = 0\)}{\(x^{2} - 10x - 24 = 0\)}{\(x^{2} + 10x + 24 = 0\)}{}{}}
\End

\Data\ID{9000069908}
\Updated{2020-07-18 09:07:04}
\Level{A}
\SubArea{Quadratic equations with complex roots}
\LANGen{
\Question{
One of the solutions of the quadratic equation 
\[ 
 2x^{2} + px + 5 = 0
\]
with a real parameter \(p\) 
is 
\[ 
 x_1 = -1 + \frac{\sqrt{6}} 
 {2} \mathrm{i}.
\]
Find the value of \(p\).}
{\(4\)}{\(- 4\)}{\(8\)}{\(- 8\)}{}{}}
\LANGpl{
\Question{
Jednym z rozwiązań równania kwadratowego 
\[ 
 2x^{2} + px + 5 = 0
\]
o rzeczywistym  parametrze \(p\) 
jest 
\[ 
 x_1 = -1 + \frac{\sqrt{6}} 
 {2} \mathrm{i}.
\]
Wyznacz wartość  \(p\).}
{\(4\)}{\(- 4\)}{\(8\)}{\(- 8\)}{}{}}
\LANGcs{
\Question{
Kvadratická rovnice s reálnými koeficienty
 \[2x^{2} + px + 5 = 0\]
a reálným parametrem \( p\) má jeden kořen 
\[x_1 = -1 + \frac{\sqrt{6}} 
 {2} \mathrm{i}.\]
Koeficient \(p\) 
má hodnotu:}
{\(4\)}{\(- 4\)}{\(8\)}{\(- 8\)}{}{}}
\LANGsk{
\Question{
Kvadratická rovnica s reálnymi koeficientmi 
\[ 
 2x^{2} + px + 5 = 0
\]
a s reálnym parametrom \(p\) 
má koreň
\[ 
 x_1 = -1 + \frac{\sqrt{6}} 
 {2} \mathrm{i}.
\]
Nájdite hodnotu \(p\).}
{\(4\)}{\(- 4\)}{\(8\)}{\(- 8\)}{}{}}
\LANGes{
\Question{
Una de las soluciones de la ecuación cuadrática
\[ 
 2x^{2} + px + 5 = 0
\]
con un parámetro real \(p\) 
es 
\[ 
 x_1 = -1 + \frac{\sqrt{6}} 
 {2} \mathrm{i}.
\]
Determina el valor de \(p\).}
{\(4\)}{\(- 4\)}{\(8\)}{\(- 8\)}{}{}}
\End

\Data\ID{9000069909}
\Updated{2020-07-18 09:07:30}
\Level{A}
\SubArea{Quadratic equations with complex roots}
\LANGen{
\Question{
One of the solutions of the quadratic equation 
\[ 
 9x^{2} - 6x + p = 0
\]
with a real parameter \(p\) 
is 
\[ 
x_1=\frac{1} 
{3} + \mathrm{i}.
\]
Find the value of \(p\).}
{\(10\)}{\(- 10\)}{\(3\)}{\(- 1\)}{}{}}
\LANGpl{
\Question{
Jednym z rozwiązań równania kwadratowego 
\[ 
 9x^{2} - 6x + p = 0
\]
o rzeczywistym  parametrze \(p\) 
jest 
\[ 
x_1=\frac{1} 
{3} + \mathrm{i}.
\]
Wyznacz wartość \(p\).}
{\(10\)}{\(- 10\)}{\(3\)}{\(- 1\)}{}{}}
\LANGcs{
\Question{
Kvadratická rovnice s reálnými koeficienty 
\[9x^{2} - 6x + p = 0\]
a reálným parametrem \(p\) má jeden kořen \[x_1=\frac{1}
{3} + \mathrm{i}.\]
Koeficient \(p\) 
má hodnotu:}
{\(10\)}{\(- 10\)}{\(3\)}{\(- 1\)}{}{}}
\LANGsk{
\Question{
Kvadratická rovnica s reálnymi koeficientmi 
\[ 
 9x^{2} - 6x + p = 0
\]
a s reálnym parametrom \(p\) 
má koreň
\[ 
x_1=\frac{1} 
{3} + \mathrm{i}.
\]
Nájdite hodnotu \(p\).}
{\(10\)}{\(- 10\)}{\(3\)}{\(- 1\)}{}{}}
\LANGes{
\Question{
Una de las soluciones de la ecuación cuadrática
\[ 
 9x^{2} - 6x + p = 0
\]
con un parámetro real \(p\) 
es 
\[ 
x_1=\frac{1} 
{3} + \mathrm{i}.
\]
Determina el valor de \(p\).}
{\(10\)}{\(- 10\)}{\(3\)}{\(- 1\)}{}{}}
\End

\Data\ID{9000035601}
\Updated{2020-07-18 02:07:40}
\Level{B}
\SubArea{Quadratic equations with complex roots}
\LANGen{
\Question{
Find the values of the parameter \(t\in \mathbb{R}\) 
which guarantee that the following quadratic equation has solutions with nonzero
imaginary part. 
\[ 
 tx^{2} - 3x + 4t = 0
\]}
{\(\left (-\infty ;-\frac{3}
{4}\right )\cup \left (\frac{3}
{4};\infty \right )\)}{\(\left (-\frac{3}
{4}; \frac{3} 
{4}\right )\)}{\(\left (\frac{3}
{4};\infty \right )\)}{\(\left \{-\frac{3}
{4}; \frac{3} 
{4}\right \}\)}{\(\mathbb{R}\setminus \left \{-\frac{3}
{4}; \frac{3} 
{4}\right \}\)}{}}
\LANGpl{
\Question{
Wskaż wartości  parametru \(t\in \mathbb{R}\) 
tak, aby  podane równanie kwadratowe miało rozwiązania o niezerowej części urojonej.
\[ 
 tx^{2} - 3x + 4t = 0
\]}
{\(\left (-\infty ;-\frac{3}
{4}\right )\cup \left (\frac{3}
{4};\infty \right )\)}{\(\left (-\frac{3}
{4}; \frac{3} 
{4}\right )\)}{\(\left (\frac{3}
{4};\infty \right )\)}{\(\left \{-\frac{3}
{4}; \frac{3} 
{4}\right \}\)}{\(\mathbb{R}\setminus \left \{-\frac{3}
{4}; \frac{3} 
{4}\right \}\)}{}}
\LANGcs{
\Question{
Najděte množinu hodnot reálného parametru
\(t\), pro které
má rovnice \(tx^{2} - 3x + 4t = 0\) 
s neznámou \(x\in \mathbb{C}\) 
imaginární kořeny, tj. komplexní kořeny s nenulovou imaginární částí.}
{\(\left (-\infty ;-\frac{3}
{4}\right )\cup \left (\frac{3}
{4};\infty \right )\)}{\(\left (-\frac{3}
{4}; \frac{3} 
{4}\right )\)}{\(\left (\frac{3}
{4};\infty \right )\)}{\(\left \{-\frac{3}
{4}; \frac{3} 
{4}\right \}\)}{\(\mathbb{R}\setminus \left \{-\frac{3}
{4}; \frac{3} 
{4}\right \}\)}{}}
\LANGsk{
\Question{
Nájdite množinu hodnôt parametra \(t\in \mathbb{R}\), pre ktoré má daná kvadratická rovnica imaginárne korene, tj. komplexné korene s nenulovú imaginárny časťou.
\[ 
 tx^{2} - 3x + 4t = 0
\]}
{\(\left (-\infty ;-\frac{3}
{4}\right )\cup \left (\frac{3}
{4};\infty \right )\)}{\(\left (-\frac{3}
{4}; \frac{3} 
{4}\right )\)}{\(\left (\frac{3}
{4};\infty \right )\)}{\(\left \{-\frac{3}
{4}; \frac{3} 
{4}\right \}\)}{\(\mathbb{R}\setminus \left \{-\frac{3}
{4}; \frac{3} 
{4}\right \}\)}{}}
\LANGes{
\Question{
Determina los valores del parámetro \(t\in \mathbb{R}\) 
suponiendo que la siguiente ecuación tiene soluciones complejas con una parte imaginaria distinta de cero. 
\[ 
 tx^{2} - 3x + 4t = 0
\]}
{\(\left (-\infty ;-\frac{3}
{4}\right )\cup \left (\frac{3}
{4};\infty \right )\)}{\(\left (-\frac{3}
{4}; \frac{3} 
{4}\right )\)}{\(\left (\frac{3}
{4};\infty \right )\)}{\(\left \{-\frac{3}
{4}; \frac{3} 
{4}\right \}\)}{\(\mathbb{R}\setminus \left \{-\frac{3}
{4}; \frac{3} 
{4}\right \}\)}{}}
\End

\Data\ID{9000035605}
\Updated{2020-07-18 02:07:07}
\Level{B}
\SubArea{Quadratic equations with complex roots}
\LANGen{
\Question{
The number \(\cos \frac{7}
{6}\pi + \mathrm{i}\sin \frac{7}
{6}\pi \) 
is a solution of a quadratic equation with real valued coefficients. Find the second
solution.}
{\(\cos \frac{5}
{6}\pi + \mathrm{i}\sin \frac{5}
{6}\pi \)}{\(\cos \frac{1}
{6}\pi + \mathrm{i}\sin \frac{1}
{6}\pi \)}{\(\cos \frac{7}
{6}\pi + \mathrm{i}\sin \frac{7}
{6}\pi \)}{\(\cos \frac{11}
 {6} \pi + \mathrm{i}\sin \frac{11}
 {6} \pi \)}{}{}}
\LANGpl{
\Question{
Liczba \(\cos \frac{7}
{6}\pi + \mathrm{i}\sin \frac{7}
{6}\pi \) 
jest rozwiązaniem równania kwadratowego o współczynnikach rzeczywistych. Wskaż drugie rozwiązanie.}
{\(\cos \frac{5}
{6}\pi + \mathrm{i}\sin \frac{5}
{6}\pi \)}{\(\cos \frac{1}
{6}\pi + \mathrm{i}\sin \frac{1}
{6}\pi \)}{\(\cos \frac{7}
{6}\pi + \mathrm{i}\sin \frac{7}
{6}\pi \)}{\(\cos \frac{11}
 {6} \pi + \mathrm{i}\sin \frac{11}
 {6} \pi \)}{}{}}
\LANGcs{
\Question{
Číslo \(\cos \frac{7}
{6}\pi + \mathrm{i}\sin \frac{7}
{6}\pi \) 
je kořenem jisté kvadratické rovnice s reálnými koeficienty. Druhý kořen
této rovnice je:}
{\(\cos \frac{5}
{6}\pi + \mathrm{i}\sin \frac{5}
{6}\pi \)}{\(\cos \frac{1}
{6}\pi + \mathrm{i}\sin \frac{1}
{6}\pi \)}{\(\cos \frac{7}
{6}\pi + \mathrm{i}\sin \frac{7}
{6}\pi \)}{\(\cos \frac{11}
 {6} \pi + \mathrm{i}\sin \frac{11}
 {6} \pi \)}{}{}}
\LANGsk{
\Question{
Číslo \(\cos \frac{7}
{6}\pi + \mathrm{i}\sin \frac{7}
{6}\pi \) 
je koreňom kvadratickej rovnice s reálnymi koeficientami. Nájdite druhý koreň tejto rovnice.}
{\(\cos \frac{5}
{6}\pi + \mathrm{i}\sin \frac{5}
{6}\pi \)}{\(\cos \frac{1}
{6}\pi + \mathrm{i}\sin \frac{1}
{6}\pi \)}{\(\cos \frac{7}
{6}\pi + \mathrm{i}\sin \frac{7}
{6}\pi \)}{\(\cos \frac{11}
 {6} \pi + \mathrm{i}\sin \frac{11}
 {6} \pi \)}{}{}}
\LANGes{
\Question{
El número \(\cos \frac{7}
{6}\pi + \mathrm{i}\sin \frac{7}
{6}\pi \) 
es una solución de una ecuación cuadrática con coeficientes reales. Determina la segunda solución.}
{\(\cos \frac{5}
{6}\pi + \mathrm{i}\sin \frac{5}
{6}\pi \)}{\(\cos \frac{1}
{6}\pi + \mathrm{i}\sin \frac{1}
{6}\pi \)}{\(\cos \frac{7}
{6}\pi + \mathrm{i}\sin \frac{7}
{6}\pi \)}{\(\cos \frac{11}
 {6} \pi + \mathrm{i}\sin \frac{11}
 {6} \pi \)}{}{}}
\End

\Data\ID{9000035606}
\Updated{2020-07-18 02:07:40}
\Level{B}
\SubArea{Quadratic equations with complex roots}
\LANGen{
\Question{
Find the quadratic equation with real valued coefficients and one of the solutions
\(x_{1} = -1 + \mathrm{i}\sqrt{3}\).}
{\(x^{2} + 2x + 4 = 0\)}{\(x^{2} - 2x + 4 = 0\)}{\(x^{2} + 2x - 4 = 0\)}{\(x^{2} - 2x - 2 = 0\)}{\(x^{2} + 2x + 2 = 0\)}{}}
\LANGpl{
\Question{
Wskaż równanie kwadratowe o  współczynnikach rzeczywistych, którego jednym z rozwiązań jest 
\(x_{1} = -1 + \mathrm{i}\sqrt{3}\).}
{\(x^{2} + 2x + 4 = 0\)}{\(x^{2} - 2x + 4 = 0\)}{\(x^{2} + 2x - 4 = 0\)}{\(x^{2} - 2x - 2 = 0\)}{\(x^{2} + 2x + 2 = 0\)}{}}
\LANGcs{
\Question{
Kvadratická rovnice s reálnými koeficienty, jejíž jeden kořen je
\(x_{1} = -1 + \mathrm{i}\sqrt{3}\), má
tvar:}
{\(x^{2} + 2x + 4 = 0\)}{\(x^{2} - 2x + 4 = 0\)}{\(x^{2} + 2x - 4 = 0\)}{\(x^{2} - 2x - 2 = 0\)}{\(x^{2} + 2x + 2 = 0\)}{}}
\LANGsk{
\Question{
Nájdite kvadratickú rovnicu s reálnymi koeficientami, ktorej jeden koreň je
\(x_{1} = -1 + \mathrm{i}\sqrt{3}\).}
{\(x^{2} + 2x + 4 = 0\)}{\(x^{2} - 2x + 4 = 0\)}{\(x^{2} + 2x - 4 = 0\)}{\(x^{2} - 2x - 2 = 0\)}{\(x^{2} + 2x + 2 = 0\)}{}}
\LANGes{
\Question{
Determina la ecuación cuadrática con coeficientes reales cuya una solución es
\(x_{1} = -1 + \mathrm{i}\sqrt{3}\).}
{\(x^{2} + 2x + 4 = 0\)}{\(x^{2} - 2x + 4 = 0\)}{\(x^{2} + 2x - 4 = 0\)}{\(x^{2} - 2x - 2 = 0\)}{\(x^{2} + 2x + 2 = 0\)}{}}
\End

\Data\ID{9000039105}
\Updated{2020-07-18 02:07:15}
\Level{B}
\SubArea{Quadratic equations with complex roots}
\LANGen{
\Question{
Find the quadratic equation with real coefficients such that one of the solutions is the complex
number \(x_{1} = 1 + 2\mathrm{i}\).}
{\(x^{2} - 2x + 5 = 0\)}{\(x^{2} - 2x + 3 = 0\)}{\(x^{2} + 2x + 5 = 0\)}{\(x^{2} + 2x - 3 = 0\)}{}{}}
\LANGpl{
\Question{
Wskaż równanie kwadratowe o współczynnikach rzeczywistych tak, aby jednym z rozwiązań była liczba zespolona
 \(x_{1} = 1 + 2\mathrm{i}\).}
{\(x^{2} - 2x + 5 = 0\)}{\(x^{2} - 2x + 3 = 0\)}{\(x^{2} + 2x + 5 = 0\)}{\(x^{2} + 2x - 3 = 0\)}{}{}}
\LANGcs{
\Question{
Najděte kvadratickou rovnici s reálnými koeficienty, jejíž jeden z kořenů je komplexní číslo \(x_{1} = 1 + 2\mathrm{i}\).}
{\(x^{2} - 2x + 5 = 0\)}{\(x^{2} - 2x + 3 = 0\)}{\(x^{2} + 2x + 5 = 0\)}{\(x^{2} + 2x - 3 = 0\)}{}{}}
\LANGsk{
\Question{
Nájdite kvadratickú rovnicu s reálnymi koeficientmi, ktorej jeden z koreňov je komplexné
číslo \(x_{1} = 1 + 2\mathrm{i}\).}
{\(x^{2} - 2x + 5 = 0\)}{\(x^{2} - 2x + 3 = 0\)}{\(x^{2} + 2x + 5 = 0\)}{\(x^{2} + 2x - 3 = 0\)}{}{}}
\LANGes{
\Question{
Determina la ecuación cuadrática con coeficientes reales suponiendo que una de las soluciones es el número complejo \(x_{1} = 1 + 2\mathrm{i}\).}
{\(x^{2} - 2x + 5 = 0\)}{\(x^{2} - 2x + 3 = 0\)}{\(x^{2} + 2x + 5 = 0\)}{\(x^{2} + 2x - 3 = 0\)}{}{}}
\End

\Data\ID{9000039106}
\Updated{2020-07-18 02:07:35}
\Level{B}
\SubArea{Quadratic equations with complex roots}
\LANGen{
\Question{
Find the value of the parameter \(a\) 
which guarantees that the quadratic equation 
\[ 
 x^{2} + 2ax + a = 0
\]
has a pair of complex conjugate solutions with a nonzero imaginary part.}
{\(a\in (0;1)\)}{\(a\in [ 0;1] \)}{\(a\in (-\infty ;0)\cup (1;\infty )\)}{Such an \(a\) 
does not exist}{}{}}
\LANGpl{
\Question{
Wyznacz wartość parametru \(a\) 
tak, aby równanie kwadratowe 
\[ 
 x^{2} + 2ax + a = 0
\] miało parę rozwiązań liczby sprzężonej o niezerowej części urojonej.}
{\(a\in (0;1)\)}{\(a\in [ 0;1] \)}{\(a\in (-\infty ;0)\cup (1;\infty )\)}{Taki \(a\) nie istnieje}{}{}}
\LANGcs{
\Question{
Najděte hodnotu parametru \(a\in \mathbb{R}\), pro který má kvadratická rovnice \[x^{2} + 2ax + a = 0\] imaginární kořeny, tj. komplexní kořeny s nenulovou imaginární částí.}
{\(a\in (0;1)\)}{\(a\in \langle 0;1\rangle \)}{\(a\in (-\infty ;0)\cup (1;\infty )\)}{Takové \(a\) neexistuje.}{}{}}
\LANGsk{
\Question{
Nájdite hodnotu parametra \(a\), pre ktorý má kvadratická rovnica
\[ 
 x^{2} + 2ax + a = 0
\]
imaginárne korene, tj. komplexné korene s nenulovú imaginárny časťou .}
{\(a\in (0;1)\)}{\(a\in [ 0;1] \)}{\(a\in (-\infty ;0)\cup (1;\infty )\)}{Takové \(a\) neexistuje}{}{}}
\LANGes{
\Question{
Determina el valor del parámetro \(a\) 
suponiendo que la ecuación cuadrática  
\[ 
 x^{2} + 2ax + a = 0
\]
tiene has a pair of complex conjugate solutions with a nonzero imaginary part.
tiene un par de soluciones conjugadas complejas con una parte imaginaria distinta de cero.}
{\(a\in (0;1)\)}{\(a\in [ 0;1] \)}{\(a\in (-\infty ;0)\cup (1;\infty )\)}{Such an \(a\) 
does not exist}{}{}}
\End

\Data\ID{9000064501}
\Updated{2020-07-18 02:07:53}
\Level{B}
\SubArea{Quadratic equations with complex roots}
\LANGen{
\Question{
Find the quadratic equation with the solution
\(x_{1, 2} =\pm 2\mathrm{i}\).}
{\(x^{2} + 4 = 0\)}{\(x^{2} - 4\mathrm{i} = 0\)}{\(x^{2} - 4 = 0\)}{\(x^{2} + 4\mathrm{i} = 0\)}{}{}}
\LANGpl{
\Question{
Wskaż równanie kwadratowe, którego rozwiązaniem jest 
\(x_{1, 2} =\pm 2\mathrm{i}\).}
{\(x^{2} + 4 = 0\)}{\(x^{2} - 4\mathrm{i} = 0\)}{\(x^{2} - 4 = 0\)}{\(x^{2} + 4\mathrm{i} = 0\)}{}{}}
\LANGcs{
\Question{
Komplexní čísla \(x_{1, 2} =\pm 2\mathrm{i}\) 
jsou kořeny kvadratické rovnice:}
{\(x^{2} + 4 = 0\)}{\(x^{2} - 4\mathrm{i} = 0\)}{\(x^{2} - 4 = 0\)}{\(x^{2} + 4\mathrm{i} = 0\)}{}{}}
\LANGsk{
\Question{
Nájdite kvadratickú rovnicu, ktorej korene sú komplexné čísla:
\(x_{1, 2} =\pm 2\mathrm{i}\).}
{\(x^{2} + 4 = 0\)}{\(x^{2} - 4\mathrm{i} = 0\)}{\(x^{2} - 4 = 0\)}{\(x^{2} + 4\mathrm{i} = 0\)}{}{}}
\LANGes{
\Question{
Determina la ecuación cuadrática con la solución
\(x_{1, 2} =\pm 2\mathrm{i}\).}
{\(x^{2} + 4 = 0\)}{\(x^{2} - 4\mathrm{i} = 0\)}{\(x^{2} - 4 = 0\)}{\(x^{2} + 4\mathrm{i} = 0\)}{}{}}
\End

\Data\ID{1003107501}
\Updated{2020-07-18 01:07:01}
\Level{C}
\SubArea{Quadratic equations with complex roots}
\LANGen{
\Question{
Find the complex roots of the following quadratic equation. 
\[ 2\mathrm{i}\,x^2 - 5\mathrm{i}\,x = 0 \]}
{\( x_1=0\text{, }\  x_2 = \frac52 \)}{\( x_1=0\text{, }\  x_2 = \frac52\mathrm{i} \)}{\( x_1=0\text{, }\  x_2 = -\frac52 \)}{\( x_1=0\text{, }\ x_2 = -\frac52\mathrm{i} \)}{}{}}
\LANGpl{
\Question{
Wyznacz pierwiastki złożone podanego równania kwadratowego.
\[ 2\mathrm{i}\,x^2 - 5\mathrm{i}\,x = 0 \]}
{\( x_1=0\text{, }\  x_2 = \frac52 \)}{\( x_1=0\text{, }\  x_2 = \frac52\mathrm{i} \)}{\( x_1=0\text{, }\  x_2 = -\frac52 \)}{\( x_1=0\text{, }\ x_2 = -\frac52\mathrm{i} \)}{}{}}
\LANGcs{
\Question{
Určete komplexní kořeny kvadratické rovnice. 
\[ 2\mathrm{i}\,x^2 - 5\mathrm{i}\,x = 0 \]}
{\( x_1=0\text{, }\  x_2 = \frac52 \)}{\( x_1=0\text{, }\  x_2 = \frac52\mathrm{i} \)}{\( x_1=0\text{, }\  x_2 = -\frac52 \)}{\( x_1=0\text{, }\ x_2 = -\frac52\mathrm{i} \)}{}{}}
\LANGsk{
\Question{
Určte komplexné korene kvadratickej rovnice.
\[ 2\mathrm{i}\,x^2 - 5\mathrm{i}\,x = 0 \]}
{\( x_1=0\text{, }\  x_2 = \frac52 \)}{\( x_1=0\text{, }\  x_2 = \frac52\mathrm{i} \)}{\( x_1=0\text{, }\  x_2 = -\frac52 \)}{\( x_1=0\text{, }\ x_2 = -\frac52\mathrm{i} \)}{}{}}
\LANGes{
\Question{
Halla el conjunto de todas las raíces complejas de la siguiente ecuación cuadrática.
\[ 2\mathrm{i}\,x^2 - 5\mathrm{i}\,x = 0 \]}
{\( x_1=0\text{, }\  x_2 = \frac52 \)}{\( x_1=0\text{, }\  x_2 = \frac52\mathrm{i} \)}{\( x_1=0\text{, }\  x_2 = -\frac52 \)}{\( x_1=0\text{, }\ x_2 = -\frac52\mathrm{i} \)}{}{}}
\End

\Data\ID{1003107502}
\Updated{2020-07-18 01:07:35}
\Level{C}
\SubArea{Quadratic equations with complex roots}
\LANGen{
\Question{
Find the solution set of the following quadratic equation in the complex plane.
\[ x^2-(2-3\mathrm{i})x = 0 \]}
{\( \{ 0; 2-3\mathrm{i} \} \)}{\( \{ 0; -2+3\mathrm{i} \} \)}{\( \{ 0; 2+3\mathrm{i} \} \)}{\( \{ 0; -2-3\mathrm{i} \} \)}{}{}}
\LANGpl{
\Question{
Wyznacz zbiór rozwiązań podanego równania kwadratowego na płaszczyźnie złożonej.
\[ x^2-(2-3\mathrm{i})x = 0 \]}
{\( \{ 0; 2-3\mathrm{i} \} \)}{\( \{ 0; -2+3\mathrm{i} \} \)}{\( \{ 0; 2+3\mathrm{i} \} \)}{\( \{ 0; -2-3\mathrm{i} \} \)}{}{}}
\LANGcs{
\Question{
Určete množinu komplexních kořenů rovnice.
\[ x^2-(2-3\mathrm{i})x = 0 \]}
{\( \{ 0; 2-3\mathrm{i} \} \)}{\( \{ 0; -2+3\mathrm{i} \} \)}{\( \{ 0; 2+3\mathrm{i} \} \)}{\( \{ 0; -2-3\mathrm{i} \} \)}{}{}}
\LANGsk{
\Question{
Určte množinu komplexných koreňov rovnice.
\[ x^2-(2-3\mathrm{i})x = 0 \]}
{\( \{ 0; 2-3\mathrm{i} \} \)}{\( \{ 0; -2+3\mathrm{i} \} \)}{\( \{ 0; 2+3\mathrm{i} \} \)}{\( \{ 0; -2-3\mathrm{i} \} \)}{}{}}
\LANGes{
\Question{
Halla el conjunto de todas las raíces complejas de la siguiente ecuación cuadrática en el plano complejo. 
\[ x^2-(2-3\mathrm{i})x = 0 \]}
{\( \{ 0; 2-3\mathrm{i} \} \)}{\( \{ 0; -2+3\mathrm{i} \} \)}{\( \{ 0; 2+3\mathrm{i} \} \)}{\( \{ 0; -2-3\mathrm{i} \} \)}{}{}}
\End

\Data\ID{1003107503}
\Updated{2020-07-18 01:07:58}
\Level{C}
\SubArea{Quadratic equations with complex roots}
\LANGen{
\Question{
Find the complex roots of the following quadratic equation.
\[ (3-\mathrm{i})x^2-(1-2\mathrm{i})x = 0 \]}
{\( x_1=0\text{, }\ x_2=\frac12-\frac12\mathrm{i} \)}{\( x_1=0\text{, }\ x_2=-\frac12+\frac12\mathrm{i} \)}{\( x_1=0\text{, }\ x_2=\frac12+\frac12\mathrm{i} \)}{\( x_1=0\text{, }\ x_2=-\frac12-\frac12\mathrm{i} \)}{}{}}
\LANGpl{
\Question{
Wyznacz pierwiastki złożone podanego równania kwadratowego. 
\[ (3-\mathrm{i})x^2-(1-2\mathrm{i})x = 0 \]}
{\( x_1=0\text{, }\ x_2=\frac12-\frac12\mathrm{i} \)}{\( x_1=0\text{, }\ x_2=-\frac12+\frac12\mathrm{i} \)}{\( x_1=0\text{, }\ x_2=\frac12+\frac12\mathrm{i} \)}{\( x_1=0\text{, }\ x_2=-\frac12-\frac12\mathrm{i} \)}{}{}}
\LANGcs{
\Question{
Určete komplexní kořeny dané kvadratické rovnice.
\[ (3-\mathrm{i})x^2-(1-2\mathrm{i})x = 0 \]}
{\( x_1=0\text{, }\ x_2=\frac12-\frac12\mathrm{i} \)}{\( x_1=0\text{, }\ x_2=-\frac12+\frac12\mathrm{i} \)}{\( x_1=0\text{, }\ x_2=\frac12+\frac12\mathrm{i} \)}{\( x_1=0\text{, }\ x_2=-\frac12-\frac12\mathrm{i} \)}{}{}}
\LANGsk{
\Question{
Určte komplexné korene danej kvadratickej rovnice.
\[ (3-\mathrm{i})x^2-(1-2\mathrm{i})x = 0 \]}
{\( x_1=0\text{, }\ x_2=\frac12-\frac12\mathrm{i} \)}{\( x_1=0\text{, }\ x_2=-\frac12+\frac12\mathrm{i} \)}{\( x_1=0\text{, }\ x_2=\frac12+\frac12\mathrm{i} \)}{\( x_1=0\text{, }\ x_2=-\frac12-\frac12\mathrm{i} \)}{}{}}
\LANGes{
\Question{
Halla las raíces complejas de la siguiente ecuación cuadrática.
\[ (3-\mathrm{i})x^2-(1-2\mathrm{i})x = 0 \]}
{\( x_1=0\text{, }\ x_2=\frac12-\frac12\mathrm{i} \)}{\( x_1=0\text{, }\ x_2=-\frac12+\frac12\mathrm{i} \)}{\( x_1=0\text{, }\ x_2=\frac12+\frac12\mathrm{i} \)}{\( x_1=0\text{, }\ x_2=-\frac12-\frac12\mathrm{i} \)}{}{}}
\End

\Data\ID{1003107504}
\Updated{2020-07-18 01:07:15}
\Level{C}
\SubArea{Quadratic equations with complex roots}
\LANGen{
\Question{
Find the complex roots of the following quadratic equation.
\[ 16\mathrm{i}x^2- 9\mathrm{i}^3 = 0 \]}
{\( x_1=\frac34\mathrm{i}\text{, }\ x_2=-\frac34\mathrm{i} \)}{\( x_1=\frac34\text{, }\ x_2=-\frac34\)}{\( x_1=\frac43\mathrm{i}\text{, }\ x_2=-\frac43\mathrm{i} \)}{\( x_1=\frac43\text{, }\ x_2=-\frac43 \)}{}{}}
\LANGpl{
\Question{
Wyznacz pierwiastki złożone podanego równania kwadratowego. 
\[ 16\mathrm{i}x^2 - 9\mathrm{i}^3 = 0 \]}
{\( x_1=\frac34\mathrm{i}\text{, }\ x_2=-\frac34\mathrm{i} \)}{\( x_1=\frac34\text{, }\ x_2=-\frac34\)}{\( x_1=\frac43\mathrm{i}\text{, }\ x_2=-\frac43\mathrm{i} \)}{\( x_1=\frac43\text{, }\ x_2=-\frac43 \)}{}{}}
\LANGcs{
\Question{
Určete komplexní kořeny dané kvadratické rovnice.
\[ 16\mathrm{i}x^2 - 9\mathrm{i}^3 = 0 \]}
{\( x_1=\frac34\mathrm{i}\text{, }\ x_2=-\frac34\mathrm{i} \)}{\( x_1=\frac34\text{, }\ x_2=-\frac34\)}{\( x_1=\frac43\mathrm{i}\text{, }\ x_2=-\frac43\mathrm{i} \)}{\( x_1=\frac43\text{, }\ x_2=-\frac43 \)}{}{}}
\LANGsk{
\Question{
Určte komplexné korene danej kvadratickej rovnice.
\[ 16\mathrm{i}x^2 - 9\mathrm{i}^3 = 0 \]}
{\( x_1=\frac34\mathrm{i}\text{, }\ x_2=-\frac34\mathrm{i} \)}{\( x_1=\frac34\text{, }\ x_2=-\frac34\)}{\( x_1=\frac43\mathrm{i}\text{, }\ x_2=-\frac43\mathrm{i} \)}{\( x_1=\frac43\text{, }\ x_2=-\frac43 \)}{}{}}
\LANGes{
\Question{
Halla las raíces complejas de la siguiente ecuación cuadrática.
\[ 16\mathrm{i}x^2- 9\mathrm{i}^3 = 0 \]}
{\( x_1=\frac34\mathrm{i}\text{, }\ x_2=-\frac34\mathrm{i} \)}{\( x_1=\frac34\text{, }\ x_2=-\frac34\)}{\( x_1=\frac43\mathrm{i}\text{, }\ x_2=-\frac43\mathrm{i} \)}{\( x_1=\frac43\text{, }\ x_2=-\frac43 \)}{}{}}
\End

\Data\ID{1003107505}
\Updated{2020-07-18 01:07:37}
\Level{C}
\SubArea{Quadratic equations with complex roots}
\LANGen{
\Question{
Find the complex roots of the following quadratic equation.
\[ 4\mathrm{i}x^2 + 1 = 0 \]}
{\( x_1=\frac{\sqrt2}4+\frac{\sqrt2}4\mathrm{i}\text{, }\ x_2=-\frac{\sqrt2}4-\frac{\sqrt2}4\mathrm{i} \)}{\( x_1=-\frac{\sqrt2}4+\frac{\sqrt2}4\mathrm{i}\text{, }\ x_2=\frac{\sqrt2}4-\frac{\sqrt2}4\mathrm{i} \)}{\( x_1=\frac12+\frac12\mathrm{i}\text{, }\ x_2=-\frac12-\frac12\mathrm{i} \)}{\( x_1=-\frac12+\frac12\mathrm{i}\text{, }\ x_2=\frac12-\frac12\mathrm{i} \)}{}{}}
\LANGpl{
\Question{
Wyznacz pierwiastki złożone podanego równania kwadratowego. 
\[ 4\mathrm{i}x^2 + 1 = 0 \]}
{\( x_1=\frac{\sqrt2}4+\frac{\sqrt2}4\mathrm{i}\text{, }\ x_2=-\frac{\sqrt2}4-\frac{\sqrt2}4\mathrm{i} \)}{\( x_1=-\frac{\sqrt2}4+\frac{\sqrt2}4\mathrm{i}\text{, }\ x_2=\frac{\sqrt2}4-\frac{\sqrt2}4\mathrm{i} \)}{\( x_1=\frac12+\frac12\mathrm{i}\text{, }\ x_2=-\frac12-\frac12\mathrm{i} \)}{\( x_1=-\frac12+\frac12\mathrm{i}\text{, }\ x_2=\frac12-\frac12\mathrm{i} \)}{}{}}
\LANGcs{
\Question{
Určete komplexní kořeny dané kvadratické rovnice.
\[ 4\mathrm{i}x^2 + 1 = 0 \]}
{\( x_1=\frac{\sqrt2}4+\frac{\sqrt2}4\mathrm{i}\text{, }\ x_2=-\frac{\sqrt2}4-\frac{\sqrt2}4\mathrm{i} \)}{\( x_1=-\frac{\sqrt2}4+\frac{\sqrt2}4\mathrm{i}\text{, }\ x_2=\frac{\sqrt2}4-\frac{\sqrt2}4\mathrm{i} \)}{\( x_1=\frac12+\frac12\mathrm{i}\text{, }\ x_2=-\frac12-\frac12\mathrm{i} \)}{\( x_1=-\frac12+\frac12\mathrm{i}\text{, }\ x_2=\frac12-\frac12\mathrm{i} \)}{}{}}
\LANGsk{
\Question{
Určte komplexné korene danej kvadratickej rovnice.
\[ 4\mathrm{i}x^2 + 1 = 0 \]}
{\( x_1=\frac{\sqrt2}4+\frac{\sqrt2}4\mathrm{i}\text{, }\ x_2=-\frac{\sqrt2}4-\frac{\sqrt2}4\mathrm{i} \)}{\( x_1=-\frac{\sqrt2}4+\frac{\sqrt2}4\mathrm{i}\text{, }\ x_2=\frac{\sqrt2}4-\frac{\sqrt2}4\mathrm{i} \)}{\( x_1=\frac12+\frac12\mathrm{i}\text{, }\ x_2=-\frac12-\frac12\mathrm{i} \)}{\( x_1=-\frac12+\frac12\mathrm{i}\text{, }\ x_2=\frac12-\frac12\mathrm{i} \)}{}{}}
\LANGes{
\Question{
Halla las raíces complejas de la siguiente ecuación cuadrática.
\[ 4\mathrm{i}x^2 + 1 = 0 \]}
{\( x_1=\frac{\sqrt2}4+\frac{\sqrt2}4\mathrm{i}\text{, }\ x_2=-\frac{\sqrt2}4-\frac{\sqrt2}4\mathrm{i} \)}{\( x_1=-\frac{\sqrt2}4+\frac{\sqrt2}4\mathrm{i}\text{, }\ x_2=\frac{\sqrt2}4-\frac{\sqrt2}4\mathrm{i} \)}{\( x_1=\frac12+\frac12\mathrm{i}\text{, }\ x_2=-\frac12-\frac12\mathrm{i} \)}{\( x_1=-\frac12+\frac12\mathrm{i}\text{, }\ x_2=\frac12-\frac12\mathrm{i} \)}{}{}}
\End

\Data\ID{1003107506}
\Updated{2020-07-18 01:07:47}
\Level{C}
\SubArea{Quadratic equations with complex roots}
\LANGen{
\Question{
Find the complex roots of the following quadratic equation.
\[ 9x^2 + 4\mathrm{i} = 0 \]}
{\( x_1=-\frac{\sqrt2}3+\frac{\sqrt2}3\mathrm{i}\text{, }\ x_2=\frac{\sqrt2}3-\frac{\sqrt2}3\mathrm{i} \)}{\( x_1=\frac23\mathrm{i}\text{, }\ x_2=-\frac23\mathrm{i} \)}{\( x_1=\frac{\sqrt2}3+\frac{\sqrt2}3\mathrm{i}\text{, }\ x_2=-\frac{\sqrt2}3-\frac{\sqrt2}3\mathrm{i} \)}{\( x_1=\frac23\text{, }\ x_2=-\frac23 \)}{}{}}
\LANGpl{
\Question{
Wyznacz pierwiastki złożone podanego równania kwadratowego. 
\[ 9x^2 + 4\mathrm{i} = 0 \]}
{\( x_1=-\frac{\sqrt2}3+\frac{\sqrt2}3\mathrm{i}\text{, }\ x_2=\frac{\sqrt2}3-\frac{\sqrt2}3\mathrm{i} \)}{\( x_1=\frac23\mathrm{i}\text{, }\ x_2=-\frac23\mathrm{i} \)}{\( x_1=\frac{\sqrt2}3+\frac{\sqrt2}3\mathrm{i}\text{, }\ x_2=-\frac{\sqrt2}3-\frac{\sqrt2}3\mathrm{i} \)}{\( x_1=\frac23\text{, }\ x_2=-\frac23 \)}{}{}}
\LANGcs{
\Question{
Určete komplexní kořeny dané kvadratické rovnice.
\[ 9x^2 + 4\mathrm{i} = 0 \]}
{\( x_1=-\frac{\sqrt2}3+\frac{\sqrt2}3\mathrm{i}\text{, }\ x_2=\frac{\sqrt2}3-\frac{\sqrt2}3\mathrm{i} \)}{\( x_1=\frac23\mathrm{i}\text{, }\ x_2=-\frac23\mathrm{i} \)}{\( x_1=\frac{\sqrt2}3+\frac{\sqrt2}3\mathrm{i}\text{, }\ x_2=-\frac{\sqrt2}3-\frac{\sqrt2}3\mathrm{i} \)}{\( x_1=\frac23\text{, }\ x_2=-\frac23 \)}{}{}}
\LANGsk{
\Question{
Určte komplexné korene danej kvadratickej rovnice.
\[ 9x^2 + 4\mathrm{i} = 0 \]}
{\( x_1=-\frac{\sqrt2}3+\frac{\sqrt2}3\mathrm{i}\text{, }\ x_2=\frac{\sqrt2}3-\frac{\sqrt2}3\mathrm{i} \)}{\( x_1=\frac23\mathrm{i}\text{, }\ x_2=-\frac23\mathrm{i} \)}{\( x_1=\frac{\sqrt2}3+\frac{\sqrt2}3\mathrm{i}\text{, }\ x_2=-\frac{\sqrt2}3-\frac{\sqrt2}3\mathrm{i} \)}{\( x_1=\frac23\text{, }\ x_2=-\frac23 \)}{}{}}
\LANGes{
\Question{
Halla las raíces complejas de la siguiente ecuación cuadrática.
\[ 9x^2 + 4\mathrm{i} = 0 \]}
{\( x_1=-\frac{\sqrt2}3+\frac{\sqrt2}3\mathrm{i}\text{, }\ x_2=\frac{\sqrt2}3-\frac{\sqrt2}3\mathrm{i} \)}{\( x_1=\frac23\mathrm{i}\text{, }\ x_2=-\frac23\mathrm{i} \)}{\( x_1=\frac{\sqrt2}3+\frac{\sqrt2}3\mathrm{i}\text{, }\ x_2=-\frac{\sqrt2}3-\frac{\sqrt2}3\mathrm{i} \)}{\( x_1=\frac23\text{, }\ x_2=-\frac23 \)}{}{}}
\End

\Data\ID{1003107507}
\Updated{2020-07-18 01:07:26}
\Level{C}
\SubArea{Quadratic equations with complex roots}
\LANGen{
\Question{
Find the solution set of the following quadratic equation in the complex plane.
\[ x^2 +3\mathrm{i}x + 10 = 0 \]}
{\( \{-5\mathrm{i}; 2\mathrm{i}\} \)}{\( \{-2\mathrm{i}; 5\mathrm{i}\} \)}{\( \{2-5\mathrm{i}; 5-2\mathrm{i}\} \)}{\( \{-2+5\mathrm{i}; -5+2\mathrm{i}\} \)}{\( \emptyset \)}{}}
\LANGpl{
\Question{
Wyznacz zbiór rozwiązań podanego równania kwadratowego na płaszczyźnie złożonej. 
\[ x^2 +3\mathrm{i}x + 10 = 0 \]}
{\( \{-5\mathrm{i}; 2\mathrm{i}\} \)}{\( \{-2\mathrm{i}; 5\mathrm{i}\} \)}{\( \{2-5\mathrm{i}; 5-2\mathrm{i}\} \)}{\( \{-2+5\mathrm{i}; -5+2\mathrm{i}\} \)}{\( \emptyset \)}{}}
\LANGcs{
\Question{
Určete množinu komplexních kořenů dané rovnice.
\[ x^2 +3\mathrm{i}x + 10 = 0 \]}
{\( \{-5\mathrm{i}; 2\mathrm{i}\} \)}{\( \{-2\mathrm{i}; 5\mathrm{i}\} \)}{\( \{2-5\mathrm{i}; 5-2\mathrm{i}\} \)}{\( \{-2+5\mathrm{i}; -5+2\mathrm{i}\} \)}{\( \emptyset \)}{}}
\LANGsk{
\Question{
Určte množinu komplexných koreňov danej rovnice.
\[ x^2 +3\mathrm{i}x + 10 = 0 \]}
{\( \{-5\mathrm{i}; 2\mathrm{i}\} \)}{\( \{-2\mathrm{i}; 5\mathrm{i}\} \)}{\( \{2-5\mathrm{i}; 5-2\mathrm{i}\} \)}{\( \{-2+5\mathrm{i}; -5+2\mathrm{i}\} \)}{\( \emptyset \)}{}}
\LANGes{
\Question{
Halla el conjunto de soluciones de la siguiente ecuación cuadrática en el plano complejo.
\[ x^2 +3\mathrm{i}x + 10 = 0 \]}
{\( \{-5\mathrm{i}; 2\mathrm{i}\} \)}{\( \{-2\mathrm{i}; 5\mathrm{i}\} \)}{\( \{2-5\mathrm{i}; 5-2\mathrm{i}\} \)}{\( \{-2+5\mathrm{i}; -5+2\mathrm{i}\} \)}{\( \emptyset \)}{}}
\End

\Data\ID{1003107508}
\Updated{2020-07-18 01:07:44}
\Level{C}
\SubArea{Quadratic equations with complex roots}
\LANGen{
\Question{
Find the complex roots of the following quadratic equation.
\[ x^2 + (2 + 2\mathrm{i})x + 2\mathrm{i} = 0 \]}
{\( x_{1,2}=-1-\mathrm{i} \)}{\( x_{1,2}=1+\mathrm{i} \)}{\( x_1=1+\mathrm{i}\text{, }\  x_2=-1-\mathrm{i} \)}{\( x_1=1-\mathrm{i}\text{, }\  x_2=-1+\mathrm{i} \)}{}{}}
\LANGpl{
\Question{
Wyznacz pierwiastki złożone podanego równania kwadratowego. 
\[ x^2 + (2 + 2\mathrm{i})x + 2\mathrm{i} = 0 \]}
{\( x_{1,2}=-1-\mathrm{i} \)}{\( x_{1,2}=1+\mathrm{i} \)}{\( x_1=1+\mathrm{i}\text{, }\  x_2=-1-\mathrm{i} \)}{\( x_1=1-\mathrm{i}\text{, }\  x_2=-1+\mathrm{i} \)}{}{}}
\LANGcs{
\Question{
Určete komplexní kořeny dané rovnice.
\[ x^2 + (2 + 2\mathrm{i})x + 2\mathrm{i} = 0 \]}
{\( x_{1,2}=-1-\mathrm{i} \)}{\( x_{1,2}=1+\mathrm{i} \)}{\( x_1=1+\mathrm{i}\text{, }\  x_2=-1-\mathrm{i} \)}{\( x_1=1-\mathrm{i}\text{, }\  x_2=-1+\mathrm{i} \)}{}{}}
\LANGsk{
\Question{
Určte komplexné korene danej rovnice.
\[ x^2 + (2 + 2\mathrm{i})x + 2\mathrm{i} = 0 \]}
{\( x_{1,2}=-1-\mathrm{i} \)}{\( x_{1,2}=1+\mathrm{i} \)}{\( x_1=1+\mathrm{i}\text{, }\  x_2=-1-\mathrm{i} \)}{\( x_1=1-\mathrm{i}\text{, }\  x_2=-1+\mathrm{i} \)}{}{}}
\LANGes{
\Question{
Halla las raíces complejas de la siguiente ecuación cuadrática.
\[ x^2 + (2 + 2\mathrm{i})x + 2\mathrm{i} = 0 \]}
{\( x_{1,2}=-1-\mathrm{i} \)}{\( x_{1,2}=1+\mathrm{i} \)}{\( x_1=1+\mathrm{i}\text{, }\  x_2=-1-\mathrm{i} \)}{\( x_1=1-\mathrm{i}\text{, }\  x_2=-1+\mathrm{i} \)}{}{}}
\End

\Data\ID{1003107509}
\Updated{2020-07-18 01:07:00}
\Level{C}
\SubArea{Quadratic equations with complex roots}
\LANGen{
\Question{
Find the solution set of the following quadratic equation in the complex plane.
\[ (3-2\mathrm{i})x^2-(2-4\mathrm{i})x + 2 = 0 \]}
{\( \left\{ 1-\mathrm{i}; \frac1{13}+\frac5{13}\mathrm{i} \right\} \)}{\( \left\{ -1+\mathrm{i}; -\frac1{13}-\frac5{13}\mathrm{i} \right\} \)}{\( \emptyset \)}{\( \left\{ -1+\mathrm{i}; \frac1{13}+\frac5{13}\mathrm{i} \right\} \)}{\( \left\{ 1-\mathrm{i}; -\frac1{13}-\frac5{13}\mathrm{i} \right\} \)}{}}
\LANGpl{
\Question{
Wyznacz zbiór rozwiązań podanego równania kwadratowego na płaszczyźnie złożonej. 
\[ (3-2\mathrm{i})x^2-(2- 4\mathrm{i})x + 2 = 0 \]}
{\( \left\{ 1-\mathrm{i}; \frac1{13}+\frac5{13}\mathrm{i} \right\} \)}{\( \left\{ -1+\mathrm{i}; -\frac1{13}-\frac5{13}\mathrm{i} \right\} \)}{\( \emptyset \)}{\( \left\{ -1+\mathrm{i}; \frac1{13}+\frac5{13}\mathrm{i} \right\} \)}{\( \left\{ 1-\mathrm{i}; -\frac1{13}-\frac5{13}\mathrm{i} \right\} \)}{}}
\LANGcs{
\Question{
Určete množinu komplexních kořenů dané rovnice.
\[ (3-2\mathrm{i})x^2-(2-4\mathrm{i})x + 2 = 0 \]}
{\( \left\{ 1-\mathrm{i}; \frac1{13}+\frac5{13}\mathrm{i} \right\} \)}{\( \left\{ -1+\mathrm{i}; -\frac1{13}-\frac5{13}\mathrm{i} \right\} \)}{\( \emptyset \)}{\( \left\{ -1+\mathrm{i}; \frac1{13}+\frac5{13}\mathrm{i} \right\} \)}{\( \left\{ 1-\mathrm{i}; -\frac1{13}-\frac5{13}\mathrm{i} \right\} \)}{}}
\LANGsk{
\Question{
Určte množinu komplexných koreňov danej rovnice.
\[ (3-2\mathrm{i})x^2-(2-4\mathrm{i})x + 2 = 0 \]}
{\( \left\{ 1-\mathrm{i}; \frac1{13}+\frac5{13}\mathrm{i} \right\} \)}{\( \left\{ -1+\mathrm{i}; -\frac1{13}-\frac5{13}\mathrm{i} \right\} \)}{\( \emptyset \)}{\( \left\{ -1+\mathrm{i}; \frac1{13}+\frac5{13}\mathrm{i} \right\} \)}{\( \left\{ 1-\mathrm{i}; -\frac1{13}-\frac5{13}\mathrm{i} \right\} \)}{}}
\LANGes{
\Question{
Halla el conjunto de soluciones de la siguiente ecuación cuadrática en el plano complejo.
\[ (3-2\mathrm{i})x^2-(2-4\mathrm{i})x + 2 = 0 \]}
{\( \left\{ 1-\mathrm{i}; \frac1{13}+\frac5{13}\mathrm{i} \right\} \)}{\( \left\{ -1+\mathrm{i}; -\frac1{13}-\frac5{13}\mathrm{i} \right\} \)}{\( \emptyset \)}{\( \left\{ -1+\mathrm{i}; \frac1{13}+\frac5{13}\mathrm{i} \right\} \)}{\( \left\{ 1-\mathrm{i}; -\frac1{13}-\frac5{13}\mathrm{i} \right\} \)}{}}
\End

\Data\ID{1003107510}
\Updated{2020-07-18 01:07:10}
\Level{C}
\SubArea{Quadratic equations with complex roots}
\LANGen{
\Question{
Find the solution set of the following quadratic equation in the complex plane.
\[ x+\frac{\mathrm{i}}2=\frac1x+\frac2{\mathrm{i}} \]}
{\( \left\{-\frac12\mathrm{i};-2\mathrm{i} \right\} \)}{\( \left\{\frac12\mathrm{i};2\mathrm{i} \right\} \)}{\( \left\{\frac12\mathrm{i};-2\mathrm{i} \right\} \)}{\( \left\{-\frac12\mathrm{i};2\mathrm{i} \right\} \)}{\( \emptyset \)}{}}
\LANGpl{
\Question{
Wyznacz zbiór rozwiązań podanego równania kwadratowego na płaszczyźnie złożonej. 
\[ x+\frac{\mathrm{i}}2=\frac1x+\frac2{\mathrm{i}} \]}
{\( \left\{-\frac12\mathrm{i};-2\mathrm{i} \right\} \)}{\( \left\{\frac12\mathrm{i};2\mathrm{i} \right\} \)}{\( \left\{\frac12\mathrm{i};-2\mathrm{i} \right\} \)}{\( \left\{-\frac12\mathrm{i};2\mathrm{i} \right\} \)}{\( \emptyset \)}{}}
\LANGcs{
\Question{
Určete množinu komplexních kořenů dané rovnice. 
\[ x+\frac{\mathrm{i}}2=\frac1x+\frac2{\mathrm{i}} \]}
{\( \left\{-\frac12\mathrm{i};-2\mathrm{i} \right\} \)}{\( \left\{\frac12\mathrm{i};2\mathrm{i} \right\} \)}{\( \left\{\frac12\mathrm{i};-2\mathrm{i} \right\} \)}{\( \left\{-\frac12\mathrm{i};2\mathrm{i} \right\} \)}{\( \emptyset \)}{}}
\LANGsk{
\Question{
Určte množinu komplexných koreňov danej rovnice.
\[ x+\frac{\mathrm{i}}2=\frac1x+\frac2{\mathrm{i}} \]}
{\( \left\{-\frac12\mathrm{i};-2\mathrm{i} \right\} \)}{\( \left\{\frac12\mathrm{i};2\mathrm{i} \right\} \)}{\( \left\{\frac12\mathrm{i};-2\mathrm{i} \right\} \)}{\( \left\{-\frac12\mathrm{i};2\mathrm{i} \right\} \)}{\( \emptyset \)}{}}
\LANGes{
\Question{
Halla el conjunto de soluciones de la siguiente ecuación cuadrática en el plano complejo.
\[ x+\frac{\mathrm{i}}2=\frac1x+\frac2{\mathrm{i}} \]}
{\( \left\{-\frac12\mathrm{i};-2\mathrm{i} \right\} \)}{\( \left\{\frac12\mathrm{i};2\mathrm{i} \right\} \)}{\( \left\{\frac12\mathrm{i};-2\mathrm{i} \right\} \)}{\( \left\{-\frac12\mathrm{i};2\mathrm{i} \right\} \)}{\( \emptyset \)}{}}
\End

\Data\ID{1003107511}
\Updated{2020-07-18 01:07:18}
\Level{C}
\SubArea{Quadratic equations with complex roots}
\LANGen{
\Question{
Find the solution set of the following quadratic equation in the complex plane.
 \[ 2x^2-x=3\mathrm{i}x+2 \]}
{\( \left\{-\frac12+\frac12\mathrm{i}; 1+\mathrm{i} \right\} \)}{\( \left\{-\frac12+\mathrm{i}; 1+\frac12\mathrm{i} \right\} \)}{\( \left\{\frac12-\frac12\mathrm{i}; -1-\mathrm{i} \right\} \)}{\( \left\{\frac12-\mathrm{i}; -1-\frac12\mathrm{i} \right\} \)}{\( \left\{-\frac12-\frac12\mathrm{i}; 1-\mathrm{i} \right\} \)}{\( \left\{\frac12+\mathrm{i}; -1+\frac12\mathrm{i} \right\} \)}}
\LANGpl{
\Question{
Wyznacz zbiór rozwiązań  podanego równania kwadratowego na płaszczyźnie złożonej.
 \[ 2x^2-x=3\mathrm{i}x+2 \]}
{\( \left\{-\frac12+\frac12\mathrm{i}; 1+\mathrm{i} \right\} \)}{\( \left\{-\frac12+\mathrm{i}; 1+\frac12\mathrm{i} \right\} \)}{\( \left\{\frac12-\frac12\mathrm{i}; -1-\mathrm{i} \right\} \)}{\( \left\{\frac12-\mathrm{i}; -1-\frac12\mathrm{i} \right\} \)}{\( \left\{-\frac12-\frac12\mathrm{i}; 1-\mathrm{i} \right\} \)}{\( \left\{\frac12+\mathrm{i}; -1+\frac12\mathrm{i} \right\} \)}}
\LANGcs{
\Question{
Určete množinu komplexních kořenů dané rovnice.
 \[ 2x^2-x=3\mathrm{i}x+2 \]}
{\( \left\{-\frac12+\frac12\mathrm{i}; 1+\mathrm{i} \right\} \)}{\( \left\{-\frac12+\mathrm{i}; 1+\frac12\mathrm{i} \right\} \)}{\( \left\{\frac12-\frac12\mathrm{i}; -1-\mathrm{i} \right\} \)}{\( \left\{\frac12-\mathrm{i}; -1-\frac12\mathrm{i} \right\} \)}{\( \left\{-\frac12-\frac12\mathrm{i}; 1-\mathrm{i} \right\} \)}{\( \left\{\frac12+\mathrm{i}; -1+\frac12\mathrm{i} \right\} \)}}
\LANGsk{
\Question{
Určte množinu komplexných koreňov danej rovnice.
 \[ 2x^2-x=3\mathrm{i}x+2 \]}
{\( \left\{-\frac12+\frac12\mathrm{i}; 1+\mathrm{i} \right\} \)}{\( \left\{-\frac12+\mathrm{i}; 1+\frac12\mathrm{i} \right\} \)}{\( \left\{\frac12-\frac12\mathrm{i}; -1-\mathrm{i} \right\} \)}{\( \left\{\frac12-\mathrm{i}; -1-\frac12\mathrm{i} \right\} \)}{\( \left\{-\frac12-\frac12\mathrm{i}; 1-\mathrm{i} \right\} \)}{\( \left\{\frac12+\mathrm{i}; -1+\frac12\mathrm{i} \right\} \)}}
\LANGes{
\Question{
Halla el conjunto de soluciones de la siguiente ecuación cuadrática en el plano complejo.\[ 2x^2-x=3\mathrm{i}x+2 \]}
{\( \left\{-\frac12+\frac12\mathrm{i}; 1+\mathrm{i} \right\} \)}{\( \left\{-\frac12+\mathrm{i}; 1+\frac12\mathrm{i} \right\} \)}{\( \left\{\frac12-\frac12\mathrm{i}; -1-\mathrm{i} \right\} \)}{\( \left\{\frac12-\mathrm{i}; -1-\frac12\mathrm{i} \right\} \)}{\( \left\{-\frac12-\frac12\mathrm{i}; 1-\mathrm{i} \right\} \)}{\( \left\{\frac12+\mathrm{i}; -1+\frac12\mathrm{i} \right\} \)}}
\End

\Data\ID{1003109401}
\Updated{2020-07-18 01:07:24}
\Level{C}
\SubArea{Quadratic equations with complex roots}
\LANGen{
\Question{
Which of the given quadratic equations has the roots \( x_1 = 2 + \mathrm{i} \) and  \( x_2 = 1 - 3\,\mathrm{i} \)?}
{\( x^2 - (3 - 2\,\mathrm{i})x + 5 - 5\,\mathrm{i} = 0 \)}{\(x^2 + (3 - 2\,\mathrm{i})x + 5 - 5\,\mathrm{i} = 0 \)}{\( x^2 - (3 + 2\,\mathrm{i})x + 5 - 5\,\mathrm{i} = 0 \)}{\( x^2 + (3 + 2\,\mathrm{i})x + 5 - 5\,\mathrm{i} = 0 \)}{}{}}
\LANGpl{
\Question{
Wybierz równanie kwadratowe, którego pierwiastkami są  \( x_1 = 2 + \mathrm{i} \) i  \( x_2 = 1 - 3\,\mathrm{i} \).}
{\( x^2 - (3 - 2\,\mathrm{i})x + 5 - 5\,\mathrm{i} = 0 \)}{\(x^2 + (3- 2\,\mathrm{i})x + 5 - 5\,\mathrm{i} = 0 \)}{\( x^2 - (3 + 2\,\mathrm{i})x + 5 - 5\,\mathrm{i} = 0 \)}{\( x^2 + (3 + 2\,\mathrm{i})x + 5 - 5\,\mathrm{i} = 0 \)}{}{}}
\LANGcs{
\Question{
Vyberte kvadratickou rovnici, jejíž kořeny jsou \( x_1 = 2 + \mathrm{i} \) a  \( x_2 = 1 - 3\,\mathrm{i} \).}
{\( x^2 - (3 - 2\,\mathrm{i})x + 5 - 5\,\mathrm{i} = 0 \)}{\(x^2 + (3 - 2\,\mathrm{i})x + 5 - 5\,\mathrm{i} = 0 \)}{\( x^2 - (3 + 2\,\mathrm{i})x + 5 - 5\,\mathrm{i} = 0 \)}{\( x^2 + (3 + 2\,\mathrm{i})x + 5 - 5\,\mathrm{i} = 0 \)}{}{}}
\LANGsk{
\Question{
Vyberte kvadratickú rovnicu, ktorej korene sú \( x_1 = 2 + \mathrm{i} \) a  \( x_2 = 1 - 3\,\mathrm{i} \).}
{\( x^2 - (3 - 2\,\mathrm{i})x + 5 - 5\,\mathrm{i} = 0 \)}{\(x^2 + (3 - 2\,\mathrm{i})x + 5 - 5\,\mathrm{i} = 0 \)}{\( x^2 - (3 + 2\,\mathrm{i})x + 5 - 5\,\mathrm{i} = 0 \)}{\( x^2 + (3 + 2\,\mathrm{i})x + 5 - 5\,\mathrm{i} = 0 \)}{}{}}
\LANGes{
\Question{
¿Cuál de las ecuaciones cuadráticas da las raíces  \( x_1 = 2 + \mathrm{i} \) y  \( x_2 = 1 - 3\,\mathrm{i} \)?}
{\( x^2 - (3 - 2\,\mathrm{i})x + 5 - 5\,\mathrm{i} = 0 \)}{\(x^2 + (3 - 2\,\mathrm{i})x + 5 - 5\,\mathrm{i} = 0 \)}{\( x^2 - (3 + 2\,\mathrm{i})x + 5 - 5\,\mathrm{i} = 0 \)}{\( x^2 + (3 + 2\,\mathrm{i})x + 5 - 5\,\mathrm{i} = 0 \)}{}{}}
\End

\Data\ID{1003109402}
\Updated{2020-07-18 01:07:21}
\Level{C}
\SubArea{Quadratic equations with complex roots}
\LANGen{
\Question{
Which of the given quadratic equations has the roots \( x_1 = 1 +\mathrm{i} \) and \( x_2 = (1 +\mathrm{i})^2 \)?}
{\( x^2 - (1 + 3\,\mathrm{i})x - 2 + 2\,\mathrm{i} = 0 \)}{\( x^2 - (1 + 3\,\mathrm{i})x - 2 - 2\,\mathrm{i} = 0 \)}{\( x^2 - (1 + 3\,\mathrm{i})x + 2 + 2\,\mathrm{i} = 0 \)}{\( x^2 - (1 + 3\,\mathrm{i})x + 2 -  2\,\mathrm{i} = 0 \)}{}{}}
\LANGpl{
\Question{
Wybierz równanie kwadratowe, którego pierwiastkami są \( x_1 = 1 +\mathrm{i} \) i \( x_2 = (1 +\mathrm{i})^2 \).}
{\( x^2 - (1 + 3\,\mathrm{i})x - 2 + 2\,\mathrm{i} = 0 \)}{\( x^2 - (1 + 3\,\mathrm{i})x - 2 - 2\,\mathrm{i} = 0 \)}{\( x^2 - (1 + 3\,\mathrm{i})x + 2 + 2\,\mathrm{i} = 0 \)}{\( x^2 - (1 + 3\,\mathrm{i})x + 2 -  2\,\mathrm{i} = 0 \)}{}{}}
\LANGcs{
\Question{
Vyberte kvadratickou rovnici, jejíž kořeny jsou \( x_1 = 1 +\mathrm{i} \) a \( x_2 = (1 +\mathrm{i})^2 \).}
{\( x^2 - (1 + 3\,\mathrm{i})x - 2 + 2\,\mathrm{i} = 0 \)}{\( x^2 - (1 + 3\,\mathrm{i})x - 2 - 2\,\mathrm{i} = 0 \)}{\( x^2 - (1 + 3\,\mathrm{i})x + 2 + 2\,\mathrm{i} = 0 \)}{\( x^2 - (1 + 3\,\mathrm{i})x + 2 -  2\,\mathrm{i} = 0 \)}{}{}}
\LANGsk{
\Question{
Vyberte kvadratickú rovnicu, ktorej korene sú \( x_1 = 1 +\mathrm{i} \) a \( x_2 = (1 +\mathrm{i})^2 \).}
{\( x^2 - (1 + 3\,\mathrm{i})x - 2 + 2\,\mathrm{i} = 0 \)}{\( x^2 - (1 + 3\,\mathrm{i})x - 2 - 2\,\mathrm{i} = 0 \)}{\( x^2 - (1 + 3\,\mathrm{i})x + 2 + 2\,\mathrm{i} = 0 \)}{\( x^2 - (1 + 3\,\mathrm{i})x + 2 -  2\,\mathrm{i} = 0 \)}{}{}}
\LANGes{
\Question{
¿Cuál de las ecuaciones cuadráticas da las raíces \( x_1 = 1 +\mathrm{i} \) y \( x_2 = (1 +\mathrm{i})^2 \)?}
{\( x^2 - (1 + 3\,\mathrm{i})x - 2 + 2\,\mathrm{i} = 0 \)}{\( x^2 - (1 + 3\,\mathrm{i})x - 2 - 2\,\mathrm{i} = 0 \)}{\( x^2 - (1 + 3\,\mathrm{i})x + 2 + 2\,\mathrm{i} = 0 \)}{\( x^2 - (1 + 3\,\mathrm{i})x + 2 -  2\,\mathrm{i} = 0 \)}{}{}}
\End

\Data\ID{1003109403}
\Updated{2020-07-18 01:07:36}
\Level{C}
\SubArea{Quadratic equations with complex roots}
\LANGen{
\Question{
One of the following equations has the solutions \( x_1=\frac12-\mathrm{i} \) and \( x_2=-\frac12+2\,\mathrm{i} \). Find this equation.}
{\( 4x^2-4\,\mathrm{i}\,x+7+6\,\mathrm{i}=0 \)}{\( 4x^2-4\,\mathrm{i}\,x-9+3\,\mathrm{i}=0 \)}{\( 4x^2+4\,\mathrm{i}\,x+7+6\,\mathrm{i}\,=0 \)}{\( 4x^2+4\,\mathrm{i}\,x-9+3\,\mathrm{i}=0 \)}{}{}}
\LANGpl{
\Question{
Wyznacz równianie o rozwiązaniach  \( x_1=\frac12-\mathrm{i} \) i \( x_2=-\frac12+2\,\mathrm{i} \).}
{\( 4x^2-4\,\mathrm{i}\,x+7+6\,\mathrm{i}=0 \)}{\( 4x^2-4\,\mathrm{i}\,x-9+3\,\mathrm{i}=0 \)}{\( 4x^2+4\,\mathrm{i}\,x+7+6\,\mathrm{i}\,=0 \)}{\( 4x^2+4\,\mathrm{i}\,x-9+3\,\mathrm{i}=0 \)}{}{}}
\LANGcs{
\Question{
Jedna z následujících rovnic má kořeny \( x_1=\frac12-\mathrm{i} \), \( x_2=-\frac12+2\,\mathrm{i} \). Najděte tuto rovnici.}
{\( 4x^2-4\,\mathrm{i}\,x+7+6\,\mathrm{i}=0 \)}{\( 4x^2-4\,\mathrm{i}\,x-9+3\,\mathrm{i}=0 \)}{\( 4x^2+4\,\mathrm{i}\,x+7+6\,\mathrm{i}\,=0 \)}{\( 4x^2+4\,\mathrm{i}\,x-9+3\,\mathrm{i}=0 \)}{}{}}
\LANGsk{
\Question{
Jedna z nasledujúcich rovníc má korene \( x_1=\frac12-\mathrm{i} \), \( x_2=-\frac12+2\,\mathrm{i} \). Nájdite tuto rovnicu.}
{\( 4x^2-4\,\mathrm{i}\,x+7+6\,\mathrm{i}=0 \)}{\( 4x^2-4\,\mathrm{i}\,x-9+3\,\mathrm{i}=0 \)}{\( 4x^2+4\,\mathrm{i}\,x+7+6\,\mathrm{i}\,=0 \)}{\( 4x^2+4\,\mathrm{i}\,x-9+3\,\mathrm{i}=0 \)}{}{}}
\LANGes{
\Question{
Una de las siguientes ecuaciones cuadráticas da las raíces  \( x_1=\frac12-\mathrm{i} \) y \( x_2=-\frac12+2\,\mathrm{i} \). Halla esta ecuación.}
{\( 4x^2-4\,\mathrm{i}\,x+7+6\,\mathrm{i}=0 \)}{\( 4x^2-4\,\mathrm{i}\,x-9+3\,\mathrm{i}=0 \)}{\( 4x^2+4\,\mathrm{i}\,x+7+6\,\mathrm{i}\,=0 \)}{\( 4x^2+4\,\mathrm{i}\,x-9+3\,\mathrm{i}=0 \)}{}{}}
\End

\Data\ID{1003109404}
\Updated{2020-07-18 01:07:09}
\Level{C}
\SubArea{Quadratic equations with complex roots}
\LANGen{
\Question{
One of the  roots of the quadratic equation \( x^2 + px + 1 - 3\,\mathrm{i} = 0 \)  with a complex parameter \( p \) is \( x_1  = -\mathrm{i} \). Choose the equivalent form of the given equation.}
{\( (x + \mathrm{i})(x -3 - \mathrm{i}) = 0 \)}{\( (x + \mathrm{i})(x - 3 +\mathrm{i}) = 0 \)}{\( (x -\mathrm{i})(x- 3-\mathrm{i}) = 0 \)}{\( (x +\mathrm{i})(x + 3 + \mathrm{i}) = 0 \)}{\( (x -\mathrm{i})(x- 3 + \mathrm{i}) = 0 \)}{\( (x -\mathrm{i})(x + 3 +\mathrm{i}) = 0 \)}}
\LANGpl{
\Question{
Jednym z pierwiastków równania kwadratowego  \( x^2 + px + 1 - 3\,\mathrm{i} = 0 \)  o parametrze zespolonym \( p \) to  \( x_1  = -\mathrm{i} \). Równanie można zapisać w postaci:}
{\( (x + \mathrm{i})(x -3 - \mathrm{i}) = 0 \)}{\( (x + \mathrm{i})(x - 3 +\mathrm{i}) = 0 \)}{\( (x -\mathrm{i})(x- 3-\mathrm{i}) = 0 \)}{\( (x +\mathrm{i})(x + 3 + \mathrm{i}) = 0 \)}{\( (x -\mathrm{i})(x- 3 + \mathrm{i}) = 0 \)}{\( (x -\mathrm{i})(x + 3 +\mathrm{i}) = 0 \)}}
\LANGcs{
\Question{
Kvadratická rovnice \( x^2 + px + 1 - 3\,\mathrm{i} = 0 \)  s komplexním parametrem \( p \) má jeden kořen \( x_1  = -\mathrm{i} \). Rovnici můžeme upravit do tvaru:}
{\( (x + \mathrm{i})(x -3 - \mathrm{i}) = 0 \)}{\( (x + \mathrm{i})(x - 3 +\mathrm{i}) = 0 \)}{\( (x -\mathrm{i})(x- 3-\mathrm{i}) = 0 \)}{\( (x +\mathrm{i})(x + 3 + \mathrm{i}) = 0 \)}{\( (x -\mathrm{i})(x- 3 + \mathrm{i}) = 0 \)}{\( (x -\mathrm{i})(x + 3 +\mathrm{i}) = 0 \)}}
\LANGsk{
\Question{
Kvadratická rovnica \( x^2 + px + 1 - 3\,\mathrm{i} = 0 \)  s komplexným parametrom \( p \) má jeden koreň \( x_1  = -\mathrm{i} \). Rovnicu je možno upraviť do tvaru:}
{\( (x + \mathrm{i})(x -3 - \mathrm{i}) = 0 \)}{\( (x + \mathrm{i})(x - 3 +\mathrm{i}) = 0 \)}{\( (x -\mathrm{i})(x- 3-\mathrm{i}) = 0 \)}{\( (x +\mathrm{i})(x + 3 + \mathrm{i}) = 0 \)}{\( (x -\mathrm{i})(x- 3 + \mathrm{i}) = 0 \)}{\( (x -\mathrm{i})(x + 3 +\mathrm{i}) = 0 \)}}
\LANGes{
\Question{
Una de las raíces de la ecuación cuadrática.\( x^2 + px + 1 - 3\,\mathrm{i} = 0 \) con un parámetro complejo \( p \) es \( x_1  = -\mathrm{i} \). Elige la forma equivalente de la ecuación dada.}
{\( (x + \mathrm{i})(x -3 - \mathrm{i}) = 0 \)}{\( (x + \mathrm{i})(x - 3 +\mathrm{i}) = 0 \)}{\( (x -\mathrm{i})(x- 3-\mathrm{i}) = 0 \)}{\( (x +\mathrm{i})(x + 3 + \mathrm{i}) = 0 \)}{\( (x -\mathrm{i})(x- 3 + \mathrm{i}) = 0 \)}{\( (x -\mathrm{i})(x + 3 +\mathrm{i}) = 0 \)}}
\End

\Data\ID{9000035602}
\Updated{2020-07-18 02:07:53}
\Level{C}
\SubArea{Quadratic equations with complex roots}
\LANGen{
\Question{
Find the values of the parameter \(m\in \mathbb{C}\) 
which guarantee that the following quadratic equation has a double solution.
\[ 
 mx^{2} - 2x - 1 + \mathrm{i} = 0
\]}
{\(m = -\frac{1}
{2} -\frac{1} 
{2}\mathrm{i}\)}{\(m = -1\)}{\(m = -1 + \mathrm{i}\)}{\(m = -\frac{1}
{2} + \frac{1} 
{2}\mathrm{i}\)}{}{}}
\LANGpl{
\Question{
Wskaż wartości parametru  \(m\in \mathbb{C}\) 
tak, aby podane równanie kwadratowe miało podwójne rozwiązanie.
\[ 
 mx^{2} - 2x - 1 + \mathrm{i} = 0
\]}
{\(m = -\frac{1}
{2} -\frac{1} 
{2}\mathrm{i}\)}{\(m = -1\)}{\(m = -1 + \mathrm{i}\)}{\(m = -\frac{1}
{2} + \frac{1} 
{2}\mathrm{i}\)}{}{}}
\LANGcs{
\Question{
Rovnice \(mx^{2} - 2x - 1 + \mathrm{i} = 0\) s
neznámou \(x\in \mathbb{C}\) 
má dvojnásobný kořen pro:}
{\(m = -\frac{1}
{2} -\frac{1} 
{2}\mathrm{i}\)}{\(m = -1\)}{\(m = -1 + \mathrm{i}\)}{\(m = -\frac{1}
{2} + \frac{1} 
{2}\mathrm{i}\)}{}{}}
\LANGsk{
\Question{
Nájdite hodnoty parametra \(m\in \mathbb{C}\), pre ktoré má daná kvadratická rovnica dvojnásobný koreň.
\[ 
 mx^{2} - 2x - 1 + \mathrm{i} = 0
\]}
{\(m = -\frac{1}
{2} -\frac{1} 
{2}\mathrm{i}\)}{\(m = -1\)}{\(m = -1 + \mathrm{i}\)}{\(m = -\frac{1}
{2} + \frac{1} 
{2}\mathrm{i}\)}{}{}}
\LANGes{
\Question{
Determina los valores del parámetro \(m\in \mathbb{C}\) 
 suponiendo que la siguiente ecuación cuadrática tiene una doble solución.
\[ 
 mx^{2} - 2x - 1 + \mathrm{i} = 0
\]}
{\(m = -\frac{1}
{2} -\frac{1} 
{2}\mathrm{i}\)}{\(m = -1\)}{\(m = -1 + \mathrm{i}\)}{\(m = -\frac{1}
{2} + \frac{1} 
{2}\mathrm{i}\)}{}{}}
\End

\Data\ID{9000035604}
\Updated{2020-07-18 02:07:08}
\Level{C}
\SubArea{Quadratic equations with complex roots}
\LANGen{
\Question{
Solve the following equation. 
\[ 
 x^{2} - 2\mathrm{i}x + 3 = 0
\]}
{\(\{ -\mathrm{i};3\mathrm{i}\}\)}{\(\{ - 4\mathrm{i};4\mathrm{i}\}\)}{\(\{1 - 2\mathrm{i};1 + 2\mathrm{i}\}\)}{\(\{ - 3\mathrm{i};\mathrm{i}\}\)}{}{}}
\LANGpl{
\Question{
Rozwiąż równanie.
\[ 
 x^{2} - 2\mathrm{i}x + 3 = 0
\]}
{\(\{ -\mathrm{i};3\mathrm{i}\}\)}{\(\{ - 4\mathrm{i};4\mathrm{i}\}\)}{\(\{1 - 2\mathrm{i};1 + 2\mathrm{i}\}\)}{\(\{ - 3\mathrm{i};\mathrm{i}\}\)}{}{}}
\LANGcs{
\Question{
Množina všech komplexních řešení rovnice
\(x^{2} - 2\mathrm{i}x + 3 = 0\) je:}
{\(\{ -\mathrm{i};3\mathrm{i}\}\)}{\(\{ - 4\mathrm{i};4\mathrm{i}\}\)}{\(\{1 - 2\mathrm{i};1 + 2\mathrm{i}\}\)}{\(\{ - 3\mathrm{i};\mathrm{i}\}\)}{}{}}
\LANGsk{
\Question{
Vyriešte danú rovnicu. 
\[ 
 x^{2} - 2\mathrm{i}x + 3 = 0
\]}
{\(\{ -\mathrm{i};3\mathrm{i}\}\)}{\(\{ - 4\mathrm{i};4\mathrm{i}\}\)}{\(\{1 - 2\mathrm{i};1 + 2\mathrm{i}\}\)}{\(\{ - 3\mathrm{i};\mathrm{i}\}\)}{}{}}
\LANGes{
\Question{
Resuelve la siguiente ecuación. 
\[ 
 x^{2} - 2\mathrm{i}x + 3 = 0
\]}
{\(\{ -\mathrm{i};3\mathrm{i}\}\)}{\(\{ - 4\mathrm{i};4\mathrm{i}\}\)}{\(\{1 - 2\mathrm{i};1 + 2\mathrm{i}\}\)}{\(\{ - 3\mathrm{i};\mathrm{i}\}\)}{}{}}
\End

\Data\ID{9000035607}
\Updated{2020-07-18 02:07:33}
\Level{C}
\SubArea{Quadratic equations with complex roots}
\LANGen{
\Question{
Identify the quadratic equation with solutions
\(x_{1} = 2\mathrm{i}\),
\(x_{2} = -\mathrm{i}\).}
{\(x^{2} -\mathrm{i}x + 2 = 0\)}{\(x^{2} + \mathrm{i}x + 2 = 0\)}{\(x^{2} + \mathrm{i}x - 2 = 0\)}{\(x^{2} -\mathrm{i}x - 2 = 0\)}{}{}}
\LANGpl{
\Question{
Wskaż równanie kwadratowe, którego rozwiązaniem jest  
\(x_{1} = 2\mathrm{i}\),
\(x_{2} = -\mathrm{i}\).}
{\(x^{2} -\mathrm{i}x + 2 = 0\)}{\(x^{2} + \mathrm{i}x + 2 = 0\)}{\(x^{2} + \mathrm{i}x - 2 = 0\)}{\(x^{2} -\mathrm{i}x - 2 = 0\)}{}{}}
\LANGcs{
\Question{
Vyberte tu z rovnic, jejímiž kořeny jsou čísla
\(x_{1} = 2\mathrm{i}\),
\(x_{2} = -\mathrm{i}\).}
{\(x^{2} -\mathrm{i}x + 2 = 0\)}{\(x^{2} + \mathrm{i}x + 2 = 0\)}{\(x^{2} + \mathrm{i}x - 2 = 0\)}{\(x^{2} -\mathrm{i}x - 2 = 0\)}{}{}}
\LANGsk{
\Question{
Určte kvadratickú rovnicu, ktorej korene sú čísla 
\(x_{1} = 2\mathrm{i}\),
\(x_{2} = -\mathrm{i}\).}
{\(x^{2} -\mathrm{i}x + 2 = 0\)}{\(x^{2} + \mathrm{i}x + 2 = 0\)}{\(x^{2} + \mathrm{i}x - 2 = 0\)}{\(x^{2} -\mathrm{i}x - 2 = 0\)}{}{}}
\LANGes{
\Question{
Identifica la ecuación cuadrática con soluciones 
\(x_{1} = 2\mathrm{i}\),
\(x_{2} = -\mathrm{i}\).}
{\(x^{2} -\mathrm{i}x + 2 = 0\)}{\(x^{2} + \mathrm{i}x + 2 = 0\)}{\(x^{2} + \mathrm{i}x - 2 = 0\)}{\(x^{2} -\mathrm{i}x - 2 = 0\)}{}{}}
\End

\Data\ID{9000035608}
\Updated{2020-07-18 02:07:56}
\Level{C}
\SubArea{Quadratic equations with complex roots}
\LANGen{
\Question{
The equation 
\[ 
 x^{2} - 2\mathrm{i}x + q = 0
\]
with a parameter \(q\in \mathbb{C}\) 
has a solution \(x_{1} = 1 + 2\mathrm{i}\). Find
the second solution \(x_{2}\) 
and the parameter \(q\).}
{\(x_{2} = -1,\ q = -1 - 2\mathrm{i}\)}{\(x_{2} = -1 - 4\mathrm{i},\ q = 9 - 6\mathrm{i}\)}{\(x_{2} = 1 - 4\mathrm{i},\ q = 7 - 4\mathrm{i}\)}{\(x_{2} = 1,\ q = -1 - 2\mathrm{i}\)}{\(x_{2} = -1,\ q = 1 + 2\mathrm{i}\)}{}}
\LANGpl{
\Question{
Równanie
\[ 
 x^{2} - 2\mathrm{i}x + q = 0
\]
o  parametrze \(q\in \mathbb{C}\) 
ma rozwiązanie  \(x_{1} = 1 + 2\mathrm{i}\). Wskaż 
drugie rozwiązanie  \(x_{2}\) 
oraz  parametr \(q\).}
{\(x_{2} = -1,\ q = -1 - 2\mathrm{i}\)}{\(x_{2} = -1 - 4\mathrm{i},\ q = 9 - 6\mathrm{i}\)}{\(x_{2} = 1 - 4\mathrm{i},\ q = 7 - 4\mathrm{i}\)}{\(x_{2} = 1,\ q = -1 - 2\mathrm{i}\)}{\(x_{2} = -1,\ q = 1 + 2\mathrm{i}\)}{}}
\LANGcs{
\Question{
Rovnice \(x^{2} - 2\mathrm{i}x + q = 0\) má
jeden kořen \(x_{1} = 1 + 2\mathrm{i}\). Najděte
druhý kořen \(x_{2}\) 
a koeficient \(q\in \mathbb{C}\).}
{\(x_{2} = -1,\ q = -1 - 2\mathrm{i}\)}{\(x_{2} = -1 - 4\mathrm{i},\ q = 9 - 6\mathrm{i}\)}{\(x_{2} = 1 - 4\mathrm{i},\ q = 7 - 4\mathrm{i}\)}{\(x_{2} = 1,\ q = -1 - 2\mathrm{i}\)}{\(x_{2} = -1,\ q = 1 + 2\mathrm{i}\)}{}}
\LANGsk{
\Question{
Rovnica 
\[ 
 x^{2} - 2\mathrm{i}x + q = 0
\]
s parametrom \(q\in \mathbb{C}\) 
má jeden koreň \(x_{1} = 1 + 2\mathrm{i}\). Nájdite druhý koreň
 \(x_{2}\) 
a parameter \(q\).}
{\(x_{2} = -1,\ q = -1 - 2\mathrm{i}\)}{\(x_{2} = -1 - 4\mathrm{i},\ q = 9 - 6\mathrm{i}\)}{\(x_{2} = 1 - 4\mathrm{i},\ q = 7 - 4\mathrm{i}\)}{\(x_{2} = 1,\ q = -1 - 2\mathrm{i}\)}{\(x_{2} = -1,\ q = 1 + 2\mathrm{i}\)}{}}
\LANGes{
\Question{
La ecuación
\[ 
 x^{2} - 2\mathrm{i}x + q = 0
\]
con un parámetro \(q\in \mathbb{C}\) 
tiene una solución \(x_{1} = 1 + 2\mathrm{i}\). Determina la segunda solución \(x_{2}\) 
y el parámetro \(q\).}
{\(x_{2} = -1,\ q = -1 - 2\mathrm{i}\)}{\(x_{2} = -1 - 4\mathrm{i},\ q = 9 - 6\mathrm{i}\)}{\(x_{2} = 1 - 4\mathrm{i},\ q = 7 - 4\mathrm{i}\)}{\(x_{2} = 1,\ q = -1 - 2\mathrm{i}\)}{\(x_{2} = -1,\ q = 1 + 2\mathrm{i}\)}{}}
\End

\Data\ID{9000035609}
\Updated{2020-07-18 02:07:02}
\Level{C}
\SubArea{Quadratic equations with complex roots}
\LANGen{
\Question{
The equation 
\[ 
 x^{2} + px - 11 = 0
\]
with a parameter \(p\in \mathbb{C}\) 
has a solution \(x_{1} = 3 -\mathrm{i}\sqrt{2}\). Find
the second solution \(x_{2}\) 
and the parameter \(p\).}
{\(x_{2} = -3 -\mathrm{i}\sqrt{2},\ p = 2\mathrm{i}\sqrt{2}\)}{\(x_{2} = 3 + \mathrm{i}\sqrt{2},\ p = 6\)}{\(x_{2} = -3 -\mathrm{i}\sqrt{2},\ p = 6\)}{\(x_{2} = 3 + \mathrm{i}\sqrt{2},\ p = -2\mathrm{i}\)}{\(x_{2} = -3 -\mathrm{i}\sqrt{2},\ p = -2\mathrm{i}\sqrt{2}\)}{}}
\LANGpl{
\Question{
Równanie
\[ 
 x^{2} + px - 11 = 0
\]
o parametrze \(p\in \mathbb{C}\) 
ma rozwiązanie  \(x_{1} = 3 -\mathrm{i}\sqrt{2}\). Wskaż drugie rozwiązanie
 \(x_{2}\) 
oraz  parametr \(p\).}
{\(x_{2} = -3 -\mathrm{i}\sqrt{2},\ p = 2\mathrm{i}\sqrt{2}\)}{\(x_{2} = 3 + \mathrm{i}\sqrt{2},\ p = 6\)}{\(x_{2} = -3 -\mathrm{i}\sqrt{2},\ p = 6\)}{\(x_{2} = 3 + \mathrm{i}\sqrt{2},\ p = -2\mathrm{i}\)}{\(x_{2} = -3 -\mathrm{i}\sqrt{2},\ p = -2\mathrm{i}\sqrt{2}\)}{}}
\LANGcs{
\Question{
Rovnice \(x^{2} + px - 11 = 0\) má
jeden kořen \(x_{1} = 3 -\mathrm{i}\sqrt{2}\). Najděte
druhý kořen \(x_{2}\) 
a koeficient \(p\in \mathbb{C}\).}
{\(x_{2} = -3 -\mathrm{i}\sqrt{2},\ p = 2\mathrm{i}\sqrt{2}\)}{\(x_{2} = 3 + \mathrm{i}\sqrt{2},\ p = 6\)}{\(x_{2} = -3 -\mathrm{i}\sqrt{2},\ p = 6\)}{\(x_{2} = 3 + \mathrm{i}\sqrt{2},\ p = -2\mathrm{i}\)}{\(x_{2} = -3 -\mathrm{i}\sqrt{2},\ p = -2\mathrm{i}\sqrt{2}\)}{}}
\LANGsk{
\Question{
Rovnica 
\[ 
 x^{2} + px - 11 = 0
\]
s parametrom \(p\in \mathbb{C}\) 
má jeden koreň \(x_{1} = 3 -\mathrm{i}\sqrt{2}\). Nájdite druhý koreň
 \(x_{2}\) 
a parameter \(p\).}
{\(x_{2} = -3 -\mathrm{i}\sqrt{2},\ p = 2\mathrm{i}\sqrt{2}\)}{\(x_{2} = 3 + \mathrm{i}\sqrt{2},\ p = 6\)}{\(x_{2} = -3 -\mathrm{i}\sqrt{2},\ p = 6\)}{\(x_{2} = 3 + \mathrm{i}\sqrt{2},\ p = -2\mathrm{i}\)}{\(x_{2} = -3 -\mathrm{i}\sqrt{2},\ p = -2\mathrm{i}\sqrt{2}\)}{}}
\LANGes{
\Question{
La ecuación
\[ 
 x^{2} + px - 11 = 0
\]
con un parámetro \(p\in \mathbb{C}\) 
tiene una solución \(x_{1} = 3 -\mathrm{i}\sqrt{2}\). Determina la segunda solución \(x_{2}\) 
y el parámetro \(p\).}
{\(x_{2} = -3 -\mathrm{i}\sqrt{2},\ p = 2\mathrm{i}\sqrt{2}\)}{\(x_{2} = 3 + \mathrm{i}\sqrt{2},\ p = 6\)}{\(x_{2} = -3 -\mathrm{i}\sqrt{2},\ p = 6\)}{\(x_{2} = 3 + \mathrm{i}\sqrt{2},\ p = -2\mathrm{i}\)}{\(x_{2} = -3 -\mathrm{i}\sqrt{2},\ p = -2\mathrm{i}\sqrt{2}\)}{}}
\End

\Data\ID{9000069910}
\Updated{2020-07-18 09:07:58}
\Level{C}
\SubArea{Quadratic equations with complex roots}
\LANGen{
\Question{
Find the values of the parameter \(p\in \mathbb{R}\) 
which guarantee that the equation 
\[ 
 x^{2} + 2px + 16 = 0
\]
has solution with a nonzero imaginary part.}
{\(p\in (-4;4)\)}{\(p\in (-\infty ;4)\)}{\(p\in (4;\infty )\)}{\(p\in \{\}\)}{}{}}
\LANGpl{
\Question{
Wyznacz wartości parametru  \(p\in \mathbb{R}\) 
tak, aby równanie 
\[ 
 x^{2} + 2px + 16 = 0
\]
miało rozwiązanie o niezerowej części urojonej.}
{\(p\in (-4;4)\)}{\(p\in (-\infty ;4)\)}{\(p\in (4;\infty )\)}{\(p\in \{\}\)}{}{}}
\LANGcs{
\Question{
Určete množinu všech hodnot parametru
\(p\in \mathbb{R}\), pro které
má rovnice \(x^{2} + 2px + 16 = 0\) 
imaginární kořeny, tj. komplexní kořeny s nenulovou imaginární částí.}
{\(p\in (-4;4)\)}{\(p\in (-\infty ;4)\)}{\(p\in (4;\infty )\)}{\(p\in \{\}\)}{}{}}
\LANGsk{
\Question{
Určte množinu všetkých hodnôt parametra \(p\in \mathbb{R}\), 
pre ktoré má rovnica  
\[ 
 x^{2} + 2px + 16 = 0
\]
imaginárne korene, tj. komplexné korene s nenulovú imaginárny časťou.}
{\(p\in (-4;4)\)}{\(p\in (-\infty ;4)\)}{\(p\in (4;\infty )\)}{\(p\in \{\}\)}{}{}}
\LANGes{
\Question{
Determina los valores del parámetro  \(p\in \mathbb{R}\) 
suponiendo que la ecuación
\[ 
 x^{2} + 2px + 16 = 0
\]
tiene la solución con una parte imaginaria distinta de cero.}
{\(p\in (-4;4)\)}{\(p\in (-\infty ;4)\)}{\(p\in (4;\infty )\)}{\(p\in \{\}\)}{}{}}
\End
